%% 2-10 単位の世界ルール：SI単位　第2部
%% 第11回（X11）：換算に使う数も単位つきの量（60 s/min、24 h/d、3600 s/h、1000 m/km）。換算は「1をかける」で、単位が約分されて消える。
%% この章の部品（名前は \bkiik…）
\makeatletter
% スピードメーター（半円）  \bkiikmeter{(中心)}{半径}{針の速さ km/h（0〜180）}
\newcommand\bkiikmeter[3]{\begin{scope}[shift={#1}]
  \fill[bkInk!85, rounded corners=2mm] (-#2-4,-6) rectangle (#2+4,#2+4);
  \fill[white] (-#2,0) arc (180:0:#2) -- cycle;
  \foreach \v in {0,10,...,180} {\pgfmathtruncatemacro\bkiik@l{mod(\v,20)==0}%
    \ifnum\bkiik@l=1 \draw[bkInk, line width=0.8pt] ({180-\v}:#2*0.97) -- ({180-\v}:#2*0.8);
      \node[font=\sffamily\bfseries\fontsize{7}{7}\selectfont, inner sep=0pt] at ({180-\v}:#2*0.64) {\v};
    \else \draw[bkInk, line width=0.5pt] ({180-\v}:#2*0.97) -- ({180-\v}:#2*0.88);\fi}
  \draw[bkRed, line width=1.4pt, line cap=round] (0,0) -- ({180-#3}:#2*0.86);
  \fill[bkInk] (0,0) circle (1.6);
  \node[font=\sffamily\bfseries\footnotesize, text=white, inner sep=0pt] at (0,-3.2) {km/h};
\end{scope}}
% 単位のふだ  \bkiiktag{(中心)}{物の名前}{数と単位}
\newcommand\bkiiktag[3]{\begin{scope}[shift={#1}]
  \filldraw[fill=white, draw=bkGreen!70!black, line width=0.7pt, rounded corners=1.5mm] (-18,-6.5) rectangle (18,6.5);
  \node[font=\sffamily\footnotesize, inner sep=0pt] at (0,2.6) {#2};
  \node[font=\sffamily\bfseries\small, text=bkGreen!50!black, inner sep=0pt] at (0,-2.8) {#3};
\end{scope}}
% 数式の中で、単位に赤い斜線を引く（約分されて消える）  \bkiikc{\mathrm{min}}
\newcommand\bkiikc[1]{{\renewcommand\CancelColor{\color{bkRed}}\cancel{#1}}}
% 文の中のことば（漢字の単位）に赤い斜線を引く  \bkiikt{\R{年}{ねん}}
\newcommand\bkiikt[1]{\tikz[baseline=(bkiik@t.base)]{\node[inner sep=0pt, outer sep=0pt] (bkiik@t) {#1};
  \draw[bkRed, line width=0.7pt] ($(bkiik@t.base west)+(-0.05em,-0.12em)$) -- ($(bkiik@t.base east)+(0.05em,0.8em)$);}}
% ことばや式を、まん中に1行で置く  \bkiikline{…}
\newcommand\bkiikline[1]{\par\addvspace{3pt}{\centering #1\par}\addvspace{3pt}}
\makeatother

\bkchapter{2}{10}{\R{単位}{たん|い}の\R{世界}{せ|かい}ルール：SI\R{単位}{たん|い}}{\bklevelrange{1}{2}}{\R{単位}{たん|い}の\R{世界}{せ|かい}ルール：SI\R{単位}{たん|い}}
\bkmichishirube{2-1,2-3,2-4,2-7}{2-13,2-15,5-1,5-5}

\begin{goal}
\begin{itemize}[label={\color{bkTheme}●}, leftmargin=1.6\zw, itemsep=0pt, topsep=1pt]
\item \bklevel{1}\enspace k・m・{\textmu} などの\R{記号}{き|ごう}を\R{読}{よ}んで、\R{単位}{たん|い}を\R{書}{か}きかえられる。
\item \bklevel{2}\enspace \qty{60}{s/min}のような「\R{単位}{たん|い}つきの1」をかけて、\R{時間}{じ|かん}や\R{速}{はや}さの\R{単位}{たん|い}を\R{書}{か}きかえられる。
\item \bklevel{2}\enspace \R{単位}{たん|い}を\R{計算}{けい|さん}して、\R{式}{しき}が\R{正}{ただ}しいかを\R{確}{たし}かめられる。
\end{itemize}
\end{goal}

\begin{hatena}
\noindent\begin{minipage}[c]{0.56\linewidth}
\R{車}{くるま}のスピードメーターは、\R{時速}{じ|そく}\qty{72}{km}をさしています。1\R{時間}{じ|かん}に\qty{72}{km}\R{進}{すす}む\R{速}{はや}さです。

50\,m\R{走}{そう}の\R{速}{はや}さは、ふつう「1\R{秒}{びょう}に\R{何}{なん}m」で\R{考}{かんが}えますね。この\R{車}{くるま}は、1\R{秒}{びょう}に\R{何}{なん}m\R{進}{すす}むのかな？　$72\div60$？　それとも \R{何}{なに}か\R{別}{べつ}の\R{計算}{けい|さん}？
\end{minipage}\hfill
\begin{minipage}[c]{0.4\linewidth}\centering
\begin{tikzpicture}[x=1mm, y=1mm]
  \bkiikmeter{(0,0)}{20}{72}
  \node[circle, fill=bkYellow, text=white, font=\sffamily\bfseries\small, inner sep=0pt, minimum size=6mm] at (20,23) {？};
  \node[anchor=north, font=\sffamily\bfseries\small, text=bkInk, inner sep=0pt] at (0,-8.5) {\R{秒速}{びょう|そく}\R{何}{なん}m？};
\end{tikzpicture}
\end{minipage}
\end{hatena}

\R{単位}{たん|い}には、\R{世界中}{せ|かい|じゅう}で\R{使}{つか}う\R{共通}{きょう|つう}のルールがあります。ルールを\R{知}{し}ると、\R{単位}{たん|い}の\R{書}{か}きかえも、\R{答}{こた}えの\R{確}{たし}かめも、すっきりできます。この\R{章}{しょう}のいちばん\R{大事}{だい|じ}な\R{道具}{どう|ぐ}は、「\R{単位}{たん|い}つきの1」です。

\clearpage
%% ---------------- やってみよう・思い出そう ----------------
\begin{yattemiyou}[\R{身}{み}のまわりの\R{単位}{たん|い}をさがそう]
\begin{enumerate}[label={\bkcircled[bkGreen]{\arabic*}}, leftmargin=2.4\zw, itemsep=1pt, topsep=0pt]
\item \R{家}{いえ}や\R{学校}{がっ|こう}で、\R{数}{かず}と\R{単位}{たん|い}が\R{書}{か}いてある\R{物}{もの}を5つ\R{以上}{い|じょう}さがそう。
\item 「\R{物}{もの}」「\R{数}{かず}と\R{単位}{たん|い}」「\R{何}{なに}の\R{量}{りょう}か」を、ノートの\R{表}{ひょう}に\R{書}{か}こう。
\item m・g・L・Wなどの\R{前}{まえ}についている k や m に\R{丸}{まる}をつけよう。\R{同}{おな}じ\R{記号}{き|ごう}は、どんな\R{意味}{い|み}かな？
\end{enumerate}
\begin{bkfigure}
\begin{tikzpicture}[x=1mm, y=1mm]
  \bkiiktag{(20,9)}{ペットボトル}{500\,{\color{bkRed}m}L}
  \bkiiktag{(61,9)}{お\R{米}{こめ}のふくろ}{5\,{\color{bkRed}k}g}
  \bkiiktag{(102,9)}{\R{電子}{でん|し}レンジ}{600\,W}
  \bkiiktag{(20,-7)}{かぜ\R{薬}{ぐすり}（1\R{回}{かい}\R{分}{ぶん}）}{300\,{\color{bkRed}m}g}
  \bkiiktag{(61,-7)}{マラソン}{42.195\,{\color{bkRed}k}m}
  \bkiiktag{(102,-7)}{かん\R{電池}{でん|ち}}{1.5\,V}
\end{tikzpicture}
\bkfigcaption{\R{見}{み}つけた\R{単位}{たん|い}の\R{例}{れい}。\R{赤}{あか}い\R{字}{じ}が、\R{単位}{たん|い}の\R{前}{まえ}についた\R{記号}{き|ごう}。}
\end{bkfigure}
\end{yattemiyou}

\begin{omoidasou}
\bkrecall{2-1}{\R{量}{りょう}＝\R{数値}{すう|ち}×\R{単位}{たん|い}。\qty{5}{m}は\qty{1}{m}の5つ\R{分}{ぶん}。\R{単位}{たん|い}を\R{変}{か}えると\R{数値}{すう|ち}が\R{変}{か}わる：\qty{5}{m}＝\qty{500}{cm}}
\bkrecall{2-3}{\R{単位}{たん|い}も\R{計算}{けい|さん}できる：m÷s＝m/s。「/」は÷のしるしで、「1\R{秒}{びょう}あたり」と\R{読}{よ}める。}
\bkrecall{2-4}{\R{同}{おな}じ\R{数}{かず}どうしのわり\R{算}{ざん}は1（$5\div5=1$）。\R{分母}{ぶん|ぼ}と\R{分子}{ぶん|し}を\R{同}{おな}じ\R{数}{かず}でわる（\R{約分}{やく|ぶん}）。わるのは\R{逆数}{ぎゃく|すう}をかけるのと\R{同}{おな}じ。}
\bkrecall{2-7}{$1000=10^3$、$\dfrac{1}{1000}=10^{-3}$。\qty{1}{L}＝\qty{1000}{cm^3}}
\end{omoidasou}

\bkfukidashi{delta}{k は「1000\R{倍}{ばい}」、m は「1000\R{分}{ぶん}の1」。\R{記号}{き|ごう}の\R{意味}{い|み}がわかると、\R{書}{か}きかえは\R{計算}{けい|さん}1つでできるよ。}

\clearpage
%% ---------------- 数学の道具 ----------------
\bksection{doguu}{\R{数学}{すう|がく}の\R{道具}{どう|ぐ}}

\bksubsection{1}{\R{世界}{せ|かい}\R{共通}{きょう|つう}の7つの\R{基本}{き|ほん}\R{単位}{たん|い}}

\leavevmode\label{bk:2-10:si}むかしは、\R{国}{くに}ごとに\R{長}{なが}さの\R{単位}{たん|い}がちがいました。\R{日本}{に|ほん}の1\R{尺}{しゃく}は\R{約}{やく}\qty{30.3}{cm}、イギリスの1フィートは\qty{30.48}{cm}。にているのに、\R{同}{おな}じではありません。これでは、\R{部品}{ぶ|ひん}を\R{外国}{がい|こく}に\R{注文}{ちゅう|もん}するとこまります。\qty{30}{cm}の\R{板}{いた}がほしいのに、\R{少}{すこ}しだけ\R{長}{なが}さのちがう\R{板}{いた}がとどくかもしれないからです。

そこで\R{今}{いま}は、\R{世界中}{せ|かい|じゅう}で\R{同}{おな}じ\R{単位}{たん|い}を\R{使}{つか}うきまりがあります。これを\textbf{SI}（\R{国際}{こく|さい}\R{単位系}{たん|い|けい}）といいます。\bkidx{えすあい}{SI（\R{国際}{こく|さい}\R{単位系}{たん|い|けい}）}もとになるのは、\R{次}{つぎ}の7つの\textbf{\R{基本}{き|ほん}\R{単位}{たん|い}}です。ほかの\R{量}{りょう}の\R{単位}{たん|い}は、この7つを\R{組}{く}み\R{合}{あ}わせて\R{作}{つく}れます。

\begin{bkcenter}
\begin{bktable}{colspec={Q[l,m,wd=38mm] Q[c,m,wd=13mm] Q[l,m,wd=22mm] Q[l,m,wd=44mm]}, rowsep=0.8pt, cells={font=\small}, row{1}={font=\sffamily\bfseries\small}}
はかる\R{量}{りょう} & \R{記号}{き|ごう} & \R{読}{よ}み\R{方}{かた} & \R{身}{み}のまわりの\R{例}{れい} \\
\R{長}{なが}さ & m & メートル & \R{身長}{しん|ちょう}\R{約}{やく}\qty{1.4}{m} \\
\R{質量}{しつ|りょう}（ふだん「\R{重}{おも}さ」とよんでいる\R{量}{りょう}） & kg & キログラム & お\R{米}{こめ}のふくろ\qty{5}{kg} \\
\R{時間}{じ|かん} & s & \R{秒}{びょう} & 50\,m\R{走}{そう}\R{約}{やく}\qty{9}{s} \\
\R{電流}{でん|りゅう} & A & アンペア & \R{家}{いえ}のブレーカー\qty{30}{A} \\
\R{温度}{おん|ど} & K & ケルビン & 0\,℃は\R{約}{やく}\qty{273}{K}（2-2） \\
\R{物質}{ぶっ|しつ}の\R{量}{りょう} & mol & モル & \R{理科}{り|か}で\R{学}{まな}ぶ \\
\R{光}{ひかり}の\R{強}{つよ}さ & cd & カンデラ & ろうそく1\R{本}{ぽん}で\R{約}{やく}\qty{1}{cd} \\
\end{bktable}
\bktabcaption{SIの7つの\R{基本}{き|ほん}\R{単位}{たん|い}}
\end{bkcenter}

\R{時間}{じ|かん}の\R{記号}{き|ごう} s は、\R{英語}{えい|ご}の second（セカンド＝\R{秒}{びょう}）の\R{頭文字}{かしら|も|じ}です。\R{記号}{き|ごう}は、\R{小文字}{こ|も|じ}と\R{大文字}{おお|も|じ}を\R{区別}{く|べつ}します。A（アンペール）や K（ケルビン）のように、\R{人}{ひと}の\R{名前}{な|まえ}からつけた\R{単位}{たん|い}は\R{大文字}{おお|も|じ}で\R{書}{か}きます。

\bksubsection{2}{かけたりわったりして、\R{単位}{たん|い}を\R{組}{く}み\R{立}{た}てる}

\leavevmode\label{bk:2-10:kumitate}\R{面積}{めん|せき}や\R{速}{はや}さの\R{単位}{たん|い}は、\R{基本}{き|ほん}\R{単位}{たん|い}をかけたりわったりして\R{作}{つく}ります（2-3）。このような\R{単位}{たん|い}を\textbf{\R{組}{く}み\R{立}{た}てた\R{単位}{たん|い}}といいます。

\begin{bkfigure}
\begin{tikzpicture}[x=1mm, y=1mm]
  % 面積：たて 1 m × よこ 1 m
  \fill[bkBlue!15] (0,0) rectangle (18,18);
  \draw[bkBlue!70!black, line width=0.8pt] (0,0) rectangle (18,18);
  \node[anchor=north, font=\sffamily\footnotesize, inner sep=0pt] at (9,-1.8) {\qty{1}{m}};
  \node[anchor=east, font=\sffamily\footnotesize, inner sep=0pt] at (-1.8,9) {\qty{1}{m}};
  \node[font=\sffamily\bfseries\small, inner sep=0pt] at (9,9) {\qty{1}{m^2}};
  \node[anchor=west, font=\small, inner sep=0pt] at (22,9) {$\mathrm{m}\times\mathrm{m}=\mathrm{m^2}$};
  % 速さ：1 秒で 5 m
  \begin{scope}[shift={(70,0)}]
    \bkcar{(0,9)}{0.42}{bkRed}
    \draw[bk arrow, bkInk!70] (0,5.5) -- (32,5.5);
    \node[anchor=north, font=\sffamily\footnotesize, inner sep=0pt] at (16,3.6) {1\R{秒}{びょう}で\qty{5}{m}};
    \node[anchor=north, font=\small, inner sep=0pt] at (16,-2.4) {$\mathrm{m}\div\mathrm{s}=\mathrm{m/s}$};
  \end{scope}
\end{tikzpicture}
\bkfigcaption{\R{長}{なが}さ×\R{長}{なが}さで\R{面積}{めん|せき}の\R{単位}{たん|い}、\R{長}{なが}さ÷\R{時間}{じ|かん}で\R{速}{はや}さの\R{単位}{たん|い}ができる。}
\end{bkfigure}

\[ \text{\R{面積}{めん|せき}}\ \ \mathrm{m}\times\mathrm{m}=\mathrm{m^2}\qquad \text{\R{体積}{たい|せき}}\ \ \mathrm{m}\times\mathrm{m}\times\mathrm{m}=\mathrm{m^3}\qquad \text{\R{速}{はや}さ}\ \ \mathrm{m}\div\mathrm{s}=\mathrm{m/s} \]
\R{力}{ちから}の\R{単位}{たん|い}N（ニュートン）も、kg・m・sを\R{組}{く}み\R{合}{あ}わせた\R{単位}{たん|い}に\R{名前}{な|まえ}をつけたものです（2-15で\R{使}{つか}います）。\R{物理}{ぶつ|り}では、「\R{重}{おも}さ」は\R{物}{もの}が\R{地球}{ち|きゅう}に\R{引}{ひ}っぱられる\R{力}{ちから}のことで、\R{単位}{たん|い}はNです。kgではかるのは\R{質量}{しつ|りょう}（\R{物}{もの}そのものの\R{量}{りょう}）で、\R{月}{つき}に\R{持}{も}っていっても\R{変}{か}わりません。

\bksubsection{3}{\R{大}{おお}きさを\R{表}{あらわ}す\R{記号}{き|ごう}：k・m・{\textmu}}

\leavevmode\label{bk:2-10:settou}kmのkは「1000\R{倍}{ばい}」、mmの\R{前}{まえ}のmは「1000\R{分}{ぶん}の1」を\R{表}{あらわ}します。\R{単位}{たん|い}の\R{前}{まえ}につけて\R{大}{おお}きさを\R{表}{あらわ}す\R{記号}{き|ごう}を、\textbf{\R{接頭語}{せっ|とう|ご}}といいます。\bkidx{せっとうご}{\R{接頭語}{せっ|とう|ご}}どの\R{単位}{たん|い}にも\R{同}{おな}じ\R{意味}{い|み}でつきます（kg、kmのk、mg、mL、mmのm）。キロはギリシャ\R{語}{ご}の「1000」、ミリはラテン\R{語}{ご}の「1000」、センチはラテン\R{語}{ご}の「100」から\R{来}{き}たことばです。

\noindent\begin{minipage}[c]{0.44\linewidth}
\R{右}{みぎ}の\R{図}{ず}のように、kのついた\R{単位}{たん|い}から、\R{何}{なに}もつかない\R{単位}{たん|い}、mのついた\R{単位}{たん|い}へと、1\R{段}{だん}\R{下}{お}りるたびに\R{数}{かず}は1000\R{倍}{ばい}になります。\R{小}{ちい}さい\R{単位}{たん|い}で\R{数}{かぞ}えると、\R{同}{おな}じ\R{量}{りょう}でも、たくさん\R{数}{かぞ}えることになるからです。
\end{minipage}\hfill
\begin{minipage}[c]{0.53\linewidth}\centering
\begin{tikzpicture}[x=1mm, y=1mm]
  \foreach \n/\xx/\t in {a/0/km, b/27/m, c/54/mm} {\node[bk node, minimum width=12mm] (l\n) at (\xx,12) {\t};}
  \foreach \n/\xx/\t in {a/0/kg, b/27/g, c/54/mg} {\node[bk node, minimum width=12mm] (w\n) at (\xx,0) {\t};}
  \foreach \r in {l,w} {\draw[bk link] (\r a) -- (\r b); \draw[bk link] (\r b) -- (\r c);}
  \node[anchor=south, font=\sffamily\bfseries\footnotesize, text=bkRed, inner sep=0pt] at (13.5,16.8) {×1000};
  \node[anchor=south, font=\sffamily\bfseries\footnotesize, text=bkRed, inner sep=0pt] at (40.5,16.8) {×1000};
  \node[anchor=north, font=\sffamily\bfseries\footnotesize, text=bkBlue, inner sep=0pt] at (13.5,-4.8) {÷1000\,（\R{反対}{はん|たい}）};
  \node[anchor=north, font=\sffamily\bfseries\footnotesize, text=bkBlue, inner sep=0pt] at (40.5,-4.8) {÷1000\,（\R{反対}{はん|たい}）};
\end{tikzpicture}
\bkfigcaption{\R{小}{ちい}さい\R{単位}{たん|い}へは×1000、\R{大}{おお}きい\R{単位}{たん|い}へは÷1000。}
\end{minipage}

\begin{bkcenter}
\begin{bktable}{colspec={Q[c,m,wd=11mm] Q[l,m,wd=19mm] Q[l,m,wd=36mm] Q[l,m,wd=50mm]}, rowsep=0.6pt, cells={font=\small}, row{1}={font=\sffamily\bfseries\small}}
\R{記号}{き|ごう} & \R{読}{よ}み\R{方}{かた} & \R{大}{おお}きさ（2-7） & \R{例}{れい} \\
G & ギガ & $10^9$（10\R{億}{おく}\R{倍}{ばい}） & \R{大}{おお}きな\R{発電所}{はつ|でん|しょ}\qty{1}{GW} \\
M & メガ & $10^6$（100\R{万}{まん}\R{倍}{ばい}） & \R{風車}{ふう|しゃ}1つの\R{発電}{はつ|でん}\R{約}{やく}\qty{2}{MW} \\
k & キロ & $10^3$（1000\R{倍}{ばい}） & \qty{1}{km}＝\qty{1000}{m} \\
c & センチ & $10^{-2}$（100\R{分}{ぶん}の1） & \qty{1}{cm}＝\qty{0.01}{m} \\
m & ミリ & $10^{-3}$（1000\R{分}{ぶん}の1） & \qty{1}{mg}＝\qty{0.001}{g} \\
{\textmu} & マイクロ & $10^{-6}$（100\R{万分}{まん|ぶん}の1） & かみの\R{毛}{け}の\R{太}{ふと}さ\R{約}{やく}\qty{70}{\micro m} \\
n & ナノ & $10^{-9}$（10\R{億分}{おく|ぶん}の1） & ウイルスの\R{大}{おお}きさ\R{約}{やく}100\,nm \\
\end{bktable}
\bktabcaption{よく\R{使}{つか}う\R{接頭語}{せっ|とう|ご}}
\end{bkcenter}



\bkkigocards{\bkkigocard{{\textmu}}{マイクロ}{100\R{万分}{まん|ぶん}の1（$10^{-6}$）を\R{表}{あらわ}す\R{接頭語}{せっ|とう|ご}。\mbox{\qty{1}{\micro m}＝\qty{0.001}{mm}}}{micro}{2-10}}
{\bkkigocard{\sffamily m/s$^2$}{メートル\R{毎秒}{まい|びょう}\R{毎秒}{まい|びょう}}{1\R{秒}{びょう}ごとに\R{速}{はや}さが\R{何}{なん}m/sふえるか}{meters per second squared}{2-10}}

\bksubsection{4}{\R{換算}{かん|さん}に\R{使}{つか}う\R{数}{かず}も、\R{単位}{たん|い}つきの\R{量}{りょう}}

\leavevmode\label{bk:2-10:ataricount}1\R{分}{ぷん}は60\R{秒}{びょう}です。これを、\R{下}{した}の\R{図}{ず}のように\R{帯}{おび}で\R{表}{あらわ}してみましょう。1\R{分}{ぷん}の\R{帯}{おび}の\R{中}{なか}に、1\R{秒}{びょう}の\R{目}{め}もりが60こならびます。

\begin{bkfigure}
\begin{tikzpicture}[x=1mm, y=1mm]
  \foreach \i/\c in {0/bkTheme, 1/bkBlue} {
    \fill[\c!15] (\i*54,8) rectangle (\i*54+54,15);
    \draw[\c!80!black, line width=0.8pt] (\i*54,8) rectangle (\i*54+54,15);
    \node[font=\sffamily\footnotesize, inner sep=0pt] at (\i*54+27,11.5) {1\R{分}{ぷん}（\qty{1}{min}）};}
  \foreach \k in {0,...,120} {\draw[bkInk, line width=0.3pt] ({\k*0.9},4) -- ({\k*0.9},{ifthenelse(mod(\k,10)==0,7,5.6)});}
  \draw[bkInk, line width=0.7pt] (0,4) -- (108,4);
  \foreach \k in {0,30,60,90,120} \node[font=\sffamily\footnotesize, anchor=north, inner sep=0pt] at ({\k*0.9},2.4) {\k};
  \node[font=\sffamily\footnotesize, anchor=west, inner sep=0pt] at (111,1) {（\R{秒}{びょう}）};
  \bkbrace{(0,16.5)}{(54,16.5)}{60\,s}{1}
  \bkbrace{(54,16.5)}{(108,16.5)}{60\,s}{1}
  \draw[{Stealth[length=2mm]}-{Stealth[length=2mm]}, bkRed, line width=0.7pt] (0,-4.6) -- (108,-4.6);
  \node[font=\sffamily\bfseries\footnotesize, text=bkRed, fill=white, inner sep=1pt] at (54,-4.6) {2\R{分}{ふん}は\R{秒}{びょう}が60こずつ2つ\R{分}{ぶん}で120\R{秒}{びょう}};
\end{tikzpicture}
\bkfigcaption{1\R{分}{ぷん}の\R{帯}{おび}の\R{中}{なか}に、1\R{秒}{びょう}の\R{目}{め}もりが60こならぶ。}
\end{bkfigure}

\R{時間}{じ|かん}の\R{単位}{たん|い}にも\R{記号}{き|ごう}があります。\R{秒}{びょう}は s（second）、\R{分}{ふん}は min（minute：ミニット）、\R{時間}{じ|かん}は h（hour：アワー）、\R{日}{にち}は d（day：デイ）です。\R{年}{ねん}は、この\R{本}{ほん}では\R{記号}{き|ごう}を\R{使}{つか}わずに「\R{年}{ねん}」と\R{書}{か}きます。

「1\R{分}{ぷん}あたり60\R{秒}{びょう}」は、2-3の\R{速}{はや}さ「1\R{秒}{びょう}あたり\qty{5}{m}」（\qty{5}{m/s}）と\R{同}{おな}じ\R{形}{かたち}の\R{量}{りょう}です。だから、\R{速}{はや}さと\R{同}{おな}じように、\R{単位}{たん|い}をつけて\R{書}{か}けます。
\bkiikline{\large \qty{60}{s/min}\quad（1\R{分}{ぷん}あたり60\R{秒}{びょう}。「60\R{秒}{びょう}\R{毎分}{まい|ふん}」と\R{読}{よ}む）}
「/」は÷のしるしなので、\qty{60}{s/min}は $\dfrac{60\,\mathrm{s}}{1\,\mathrm{min}}$ とも\R{書}{か}けます。\R{上}{うえ}の\R{図}{ず}で、1\R{分}{ぷん}の\R{帯}{おび}と、60\R{秒}{びょう}\R{分}{ぶん}の\R{目}{め}もりは、\R{同}{おな}じ\R{長}{なが}さでした。\R{同}{おな}じ\R{大}{おお}きさどうしのわり\R{算}{ざん}は1です（$5\div5=1$ と\R{同}{おな}じ。2-4）。だから、
\[ \frac{60\,\mathrm{s}}{1\,\mathrm{min}}=1\qquad\qquad \frac{1\,\mathrm{min}}{60\,\mathrm{s}}=1 \]
\R{右}{みぎ}の\R{式}{しき}のように、\R{上}{うえ}と\R{下}{した}を\R{入}{い}れかえても1です。
\R{数}{かず}だけを\R{見}{み}ると $60\div1=60$ なのに「1」になるのは、\R{単位}{たん|い}もいっしょに\R{見}{み}ているからです。60\R{秒}{びょう}と1\R{分}{ぷん}は、\R{同}{おな}じ\R{長}{なが}さの\R{時間}{じ|かん}を、ちがう\R{単位}{たん|い}で\R{言}{い}っているだけです。このように、\R{上}{うえ}と\R{下}{した}が\R{同}{おな}じ\R{大}{おお}きさの\R{量}{りょう}になっている\R{分数}{ぶん|すう}を、この\R{本}{ほん}では\textbf{\R{単位}{たん|い}つきの1}とよびます。

\R{長}{なが}さでも\R{同}{おな}じです。\qty{1}{km}は\qty{1000}{m}なので、「1\,kmあたり1000\,m」の\qty{1000}{m/km}も\R{単位}{たん|い}つきの1です。

\bksubsection{5}{\R{換算}{かん|さん}は「1をかける」：\R{単位}{たん|い}が\R{約分}{やく|ぶん}されて\R{消}{き}える}

\leavevmode\label{bk:2-10:kansan}どんな\R{数}{かず}に1をかけても、\R{大}{おお}きさは\R{変}{か}わりません（$7\times1=7$）。だから、\R{量}{りょう}に「\R{単位}{たん|い}つきの1」をかけても、\R{量}{りょう}は\R{変}{か}わりません。\R{変}{か}わるのは\R{単位}{たん|い}だけです。このように、\R{量}{りょう}を\R{変}{か}えずに\R{単位}{たん|い}だけを\R{取}{と}りかえることを\textbf{\R{換算}{かん|さん}}といいます。\bkidx{かんさん}{\R{換算}{かん|さん}}

2\R{分}{ふん}を\R{秒}{びょう}に\R{換算}{かん|さん}してみましょう。\qty{2}{min}に\qty{60}{s/min}をかけます。
\[ \qty{2}{min}\times\qty{60}{s/min}=2\,\bkiikc{\mathrm{min}}\times\frac{60\,\mathrm{s}}{1\,\bkiikc{\mathrm{min}}}=2\times60\,\mathrm{s}=120\,\mathrm{s} \]
\R{分子}{ぶん|し}にも\R{分母}{ぶん|ぼ}にもminがあるので、minどうしが\R{約分}{やく|ぶん}されて\R{消}{き}え、sだけが\R{残}{のこ}りました。\R{数}{かず}の\R{約分}{やく|ぶん}（2-4）と\R{同}{おな}じしくみです。

\begin{bkcenter}
\begin{bktable}{colspec={Q[c,m,wd=50mm] Q[c,m,wd=66mm]}, rowsep=3pt, cells={font=\small}, row{1}={font=\sffamily\bfseries\small, bg=bkTheme!12}}
\R{数}{かず}の\R{約分}{やく|ぶん}（2-4） & \R{単位}{たん|い}の\R{約分}{やく|ぶん} \\
$\dfrac{3\times\bkiikc{5}}{\bkiikc{5}}=3$ & $2\,\bkiikc{\mathrm{min}}\times\dfrac{60\,\mathrm{s}}{1\,\bkiikc{\mathrm{min}}}=120\,\mathrm{s}$ \\
$5\div5=1$ なので\R{消}{け}してよい & min÷min＝1 なので\R{消}{け}してよい \\
\end{bktable}
\bktabcaption{\R{数}{かず}も\R{単位}{たん|い}も、\R{同}{おな}じものどうしは\R{約分}{やく|ぶん}できる}
\end{bkcenter}

\R{前}{まえ}のページの\R{帯}{おび}の\R{図}{ず}でも、2\R{分}{ふん}は「60\R{秒}{びょう}のまとまりが2つ」で120\R{秒}{びょう}でした。\R{単位}{たん|い}の\R{計算}{けい|さん}と、\R{図}{ず}で\R{数}{かぞ}えた\R{答}{こた}えが\R{合}{あ}っています。

\textbf{どちらの1をかける？}\quad \R{単位}{たん|い}つきの1には、$\dfrac{60\,\mathrm{s}}{1\,\mathrm{min}}$ と、\R{上下}{じょう|げ}を\R{入}{い}れかえた $\dfrac{1\,\mathrm{min}}{60\,\mathrm{s}}$ の2つがあります。\textbf{\R{消}{け}したい\R{単位}{たん|い}が\R{分母}{ぶん|ぼ}にくるほう}を\R{選}{えら}びます。もし\R{反対}{はん|たい}のほうをかけると、
\[ 2\,\mathrm{min}\times\frac{1\,\mathrm{min}}{60\,\mathrm{s}}=\frac{2}{60}\,\frac{\mathrm{min}\times\mathrm{min}}{\mathrm{s}} \]
となり、minが\R{消}{き}えずに2つにふえてしまいます。\R{時間}{じ|かん}の\R{単位}{たん|い}（sだけ）にならないので、まちがいに\R{気}{き}づけます。\R{単位}{たん|い}をつけて\R{計算}{けい|さん}すると、\R{単位}{たん|い}がまちがいを\R{教}{おし}えてくれるのです。

\R{反対}{はん|たい}に、\R{秒}{びょう}を\R{分}{ふん}に\R{直}{なお}すときは、sを\R{消}{け}したいので、sが\R{分母}{ぶん|ぼ}にくる $\dfrac{1\,\mathrm{min}}{60\,\mathrm{s}}$ をかけます。これは「\qty{60}{s/min}でわる」のと\R{同}{おな}じです（わるのは\R{逆数}{ぎゃく|すう}をかけるのと\R{同}{おな}じ。2-4）。
\[ \qty{180}{s}\div\qty{60}{s/min}=180\,\bkiikc{\mathrm{s}}\times\frac{1\,\mathrm{min}}{60\,\bkiikc{\mathrm{s}}}=\frac{180}{60}\,\mathrm{min}=3\,\mathrm{min} \]
\R{長}{なが}さも\R{同}{おな}じです。\qty{3}{km}をmにするには、kmを\R{消}{け}したいので、kmが\R{分母}{ぶん|ぼ}にくる\qty{1000}{m/km}をかけます。
\[ \qty{3}{km}\times\qty{1000}{m/km}=3\,\bkiikc{\mathrm{km}}\times\frac{1000\,\mathrm{m}}{1\,\bkiikc{\mathrm{km}}}=3000\,\mathrm{m} \]

\begin{bkcenter}
{\setlength\fboxsep{4pt}\fcolorbox{bkTheme}{white}{\begin{minipage}{0.86\linewidth}
\textbf{\R{換算}{かん|さん}の\R{手順}{て|じゅん}}\par
\noindent\hangindent=1.6\zw\hangafter=1 \makebox[1.6\zw][l]{\bkcircled[bkTheme]{1}}\R{今}{いま}の\R{単位}{たん|い}と、ほしい\R{単位}{たん|い}をつなぐ「\R{単位}{たん|い}つきの1」を\R{見}{み}つける（\qty{60}{s/min}など）。\par
\noindent\hangindent=1.6\zw\hangafter=1 \makebox[1.6\zw][l]{\bkcircled[bkTheme]{2}}\R{消}{け}したい\R{単位}{たん|い}が\R{分母}{ぶん|ぼ}にくる\R{向}{む}きにして、かける。\par
\noindent\hangindent=1.6\zw\hangafter=1 \makebox[1.6\zw][l]{\bkcircled[bkTheme]{3}}\R{単位}{たん|い}を\R{約分}{やく|ぶん}して、ほしい\R{単位}{たん|い}だけが\R{残}{のこ}ったか\R{確}{たし}かめる。
\end{minipage}}}
\end{bkcenter}

\bksubsection{6}{\R{秒}{びょう}・\R{分}{ふん}・\R{時間}{じ|かん}・\R{日}{にち}・\R{年}{ねん}をつなぐ}

\leavevmode\label{bk:2-10:jikan}\R{時間}{じ|かん}の\R{単位}{たん|い}は、\R{次}{つぎ}の「\R{単位}{たん|い}つきの1」でつながっています。\R{大}{おお}きい\R{単位}{たん|い}から\R{小}{ちい}さい\R{単位}{たん|い}へは、\R{右}{みぎ}へ\R{進}{すす}むたびにかけていきます。

\begin{bkfigure}
\begin{tikzpicture}[x=1mm, y=1mm]
  \foreach \n/\xx/\t in {y/0/{\R{年}{ねん}\\\ }, d/28/{\R{日}{にち}\\d}, h/56/{\R{時間}{じ|かん}\\h}, m/84/{\R{分}{ふん}\\min}, s/112/{\R{秒}{びょう}\\s}}
    {\node[bk node, minimum width=15mm, minimum height=11mm] (\n) at (\xx,0) {\t};}
  \draw[bk link] (y) -- (d); \draw[bk link] (d) -- (h); \draw[bk link] (h) -- (m); \draw[bk link] (m) -- (s);
  \foreach \xx/\a/\b in {14/×365/{\R{日}{にち}/\R{年}{ねん}}, 42/×24/{h/d}, 70/×60/{min/h}, 98/×60/{s/min}}
    {\node[anchor=south, font=\sffamily\bfseries\footnotesize, text=bkRed, align=center, inner sep=0pt] at (\xx,8.6) {\a\\\b};}
\end{tikzpicture}
\bkfigcaption{\R{右}{みぎ}へ\R{進}{すす}むときは、\R{赤}{あか}い「\R{単位}{たん|い}つきの1」をかける。\R{左}{ひだり}へもどるときは、わる。}
\end{bkfigure}

\begin{bkcenter}
\begin{bktable}{colspec={Q[l,m,wd=27mm] Q[l,m,wd=34mm] Q[l,m,wd=56mm]}, rowsep=1pt, cells={font=\small}, row{1}={font=\sffamily\bfseries\small, bg=bkTheme!12}}
\R{単位}{たん|い}つきの1 & \R{意味}{い|み} & なぜその\R{数}{かず}？ \\
\qty{60}{s/min} & 1\R{分}{ぷん}あたり60\R{秒}{びょう} & 1\R{分}{ぷん}を60\R{秒}{びょう}と\R{決}{き}めた（60は2・3・4・5・6でわり\R{切}{き}れて\R{便利}{べん|り}） \\
\qty{60}{min/h} & 1\R{時間}{じ|かん}あたり60\R{分}{ぷん} & 1\R{時間}{じ|かん}を60\R{分}{ぷん}と\R{決}{き}めた \\
\qty{24}{h/d} & 1\R{日}{にち}あたり24\R{時間}{じ|かん} & \R{地球}{ち|きゅう}が1\R{回}{かい}まわる1\R{日}{にち}を、24に\R{分}{わ}けると\R{決}{き}めた \\
365\,\R{日}{にち}/\R{年}{ねん} & 1\R{年}{ねん}あたり365\R{日}{にち} & \R{地球}{ち|きゅう}が\R{太陽}{たい|よう}のまわりを1\R{周}{しゅう}するのに\R{約}{やく}365\R{日}{にち}かかる（うるう\R{年}{どし}は366\R{日}{にち}） \\
\qty{3600}{s/h} & 1\R{時間}{じ|かん}あたり3600\R{秒}{びょう} & \R{下}{した}の\R{計算}{けい|さん}で\R{作}{つく}れる（$60\times60$） \\
\end{bktable}
\bktabcaption{\R{時間}{じ|かん}の\R{換算}{かん|さん}に\R{使}{つか}う「\R{単位}{たん|い}つきの1」}
\end{bkcenter}

\R{表}{ひょう}の\qty{3600}{s/h}は、\R{覚}{おぼ}えなくても\R{作}{つく}れます。1\R{時間}{じ|かん}を、まず\R{分}{ふん}に、\R{次}{つぎ}に\R{秒}{びょう}に\R{換算}{かん|さん}します。1をかけるたびに、1つ\R{前}{まえ}の\R{単位}{たん|い}が\R{約分}{やく|ぶん}されて\R{消}{き}えます。
\[ 1\,\bkiikc{\mathrm{h}}\times\frac{60\,\bkiikc{\mathrm{min}}}{1\,\bkiikc{\mathrm{h}}}\times\frac{60\,\mathrm{s}}{1\,\bkiikc{\mathrm{min}}}=1\times60\times60\,\mathrm{s}=3600\,\mathrm{s} \]
だから\qty{1}{h}＝\qty{3600}{s}で、「1\R{時間}{じ|かん}あたり3600\R{秒}{びょう}」の\qty{3600}{s/h}も\R{単位}{たん|い}つきの1です。

\R{年}{ねん}のように\R{日本語}{に|ほん|ご}で\R{書}{か}いた\R{単位}{たん|い}も、\R{同}{おな}じように\R{約分}{やく|ぶん}できます。2\R{年}{ねん}は\R{何日}{なん|にち}かな？
\bkiikline{2\,\bkiikt{\R{年}{ねん}}×365\,\R{日}{にち}/\bkiikt{\R{年}{ねん}}＝$2\times365$\,\R{日}{にち}＝730\,\R{日}{にち}\quad（うるう\R{年}{どし}がないとき）}

\bksubsection{7}{\R{時速}{じ|そく}を\R{秒速}{びょう|そく}に\R{直}{なお}す}

\leavevmode\label{bk:2-10:jisoku}\bkhatenaref の\R{時速}{じ|そく}\qty{72}{km}は、\qty{72}{km/h}と\R{書}{か}き、「72キロメートル\R{毎時}{まい|じ}」と\R{読}{よ}みます。1\R{時間}{じ|かん}あたり\qty{72}{km}\R{進}{すす}む\R{速}{はや}さです。1\R{秒}{びょう}あたり\R{何}{なん}m\R{進}{すす}むかを、2つの\R{方法}{ほう|ほう}で\R{求}{もと}めてみましょう。

\textbf{\R{方法}{ほう|ほう}A：\R{図}{ず}で\R{考}{かんが}える}\quad 1\R{時間}{じ|かん}は3600\R{秒}{びょう}、\qty{72}{km}は\qty{72000}{m}です（\qty{1000}{m/km}をかけた）。1\R{時間}{じ|かん}に\R{進}{すす}む\qty{72000}{m}を、3600\R{秒}{びょう}\R{分}{ぶん}に\R{等}{ひと}しく\R{分}{わ}けると、その1つ\R{分}{ぶん}が1\R{秒}{びょう}に\R{進}{すす}む\R{道}{みち}のりです（2-3の「1あたりの\R{量}{りょう}」）。

\begin{bkfigure}
\begin{tikzpicture}[x=1mm, y=1mm]
  \fill[bkBlue!14] (0,0) rectangle (118,8);
  \fill[bkOrange!55] (0,0) rectangle (6,8);
  \foreach \k in {1,...,7} {\draw[bkBlue!70!black, line width=0.4pt] (\k*6,0) -- (\k*6,8);}
  \foreach \k in {0,1,2} {\draw[bkBlue!70!black, line width=0.4pt] (100+\k*6,0) -- (100+\k*6,8);}
  \fill[white] (50,-0.5) rectangle (92,8.5);
  \node[font=\sffamily\bfseries\large, text=bkBlue!70!black, inner sep=0pt] at (71,4) {…};
  \draw[bkBlue!70!black, line width=0.8pt] (0,0) -- (50,0) (0,8) -- (50,8) (92,0) -- (118,0) (92,8) -- (118,8) (0,0) -- (0,8) (118,0) -- (118,8);
  \bkbrace{(0,9.5)}{(118,9.5)}{1\R{時間}{じ|かん}（3600\R{秒}{びょう}）で\qty{72000}{m}}{1}
  \node[anchor=north, font=\sffamily\footnotesize, inner sep=0pt] at (71,-2) {3600こに\R{等}{ひと}しく\R{分}{わ}ける};
  \draw[bkOrange!80!black, line width=0.7pt] (3,-0.6) -- (3,-6.5);
  \node[anchor=north west, font=\sffamily\bfseries\footnotesize, text=bkOrange!80!black, inner sep=0pt] at (0,-7.5) {1\R{秒}{びょう}\R{分}{ぶん}＝$72000\div3600=20$（m）};
\end{tikzpicture}
\bkfigcaption{1\R{時間}{じ|かん}の\R{道}{みち}のりを3600に\R{分}{わ}けた1つ\R{分}{ぶん}が、1\R{秒}{びょう}の\R{道}{みち}のり。}
\end{bkfigure}

\textbf{\R{方法}{ほう|ほう}B：\R{単位}{たん|い}つきの1をかける}\quad kmをmにしたいので\qty{1000}{m/km}を、hをsにしたいので\qty{3600}{s/h}を\R{使}{つか}います。\qty{72}{km/h}では、hは\R{分母}{ぶん|ぼ}にあります。\R{分母}{ぶん|ぼ}のhを\R{消}{け}すには、hが\R{分子}{ぶん|し}にくる $\dfrac{1\,\mathrm{h}}{3600\,\mathrm{s}}$ をかけます（\qty{3600}{s/h}でわるのと\R{同}{おな}じ）。
\begin{align*}
\qty{72}{km/h}&=72\,\frac{\bkiikc{\mathrm{km}}}{\bkiikc{\mathrm{h}}}\times\frac{1000\,\mathrm{m}}{1\,\bkiikc{\mathrm{km}}}\times\frac{1\,\bkiikc{\mathrm{h}}}{3600\,\mathrm{s}}\\
&=\frac{72\times1000}{3600}\,\frac{\mathrm{m}}{\mathrm{s}}=\frac{72000}{3600}\,\mathrm{m/s}=20\,\mathrm{m/s}
\end{align*}
kmもhも\R{約分}{やく|ぶん}されて\R{消}{き}え、m/sが\R{残}{のこ}りました。AでもBでも\R{秒速}{びょう|そく}\qty{20}{m}です。これが\bkhatenaref の\R{答}{こた}えです。

\textbf{$72\div60$ だと、どうなる？}\quad 60は「1\R{時間}{じ|かん}あたり60\R{分}{ぷん}」の\qty{60}{min/h}です。\R{単位}{たん|い}もつけて\R{計算}{けい|さん}すると、
\[ \qty{72}{km/h}\div\qty{60}{min/h}=72\,\frac{\mathrm{km}}{\bkiikc{\mathrm{h}}}\times\frac{1\,\bkiikc{\mathrm{h}}}{60\,\mathrm{min}}=1.2\,\mathrm{km/min} \]
これは「1\R{分}{ぷん}あたり\qty{1.2}{km}」で、まちがいではありませんが、ほしかった「1\R{秒}{びょう}あたり\R{何}{なん}m」ではありません。\R{単位}{たん|い}を\R{見}{み}れば、まだkmと\R{分}{ふん}のままだとわかります。

\textbf{\R{近道}{ちか|みち}：÷3.6}\quad \R{方法}{ほう|ほう}Bでは、いつも $\times1000\div3600$ をしています。ここで、3600は1000の3.6\R{倍}{ばい}です（$3600\div1000=3.6$）。だから、$\times1000\div3600$ は「÷3.6」と\R{同}{おな}じです。だから、\textbf{\R{時速}{じ|そく}を\R{秒速}{びょう|そく}にするには÷3.6}、\R{秒速}{びょう|そく}を\R{時速}{じ|そく}にするには×3.6です。\R{下}{した}の\R{表}{ひょう}で\R{確}{たし}かめましょう。
\begin{bkcenter}
\begin{tblr}{colspec={Q[l,m] *{3}{Q[c,m,wd=9mm]}}, rowsep=0.5pt, colsep=3pt, hlines={0.5pt,bkLine}, vlines={0.5pt,bkLine}, cells={font=\small}, column{1}={bg=bkTheme!12}}
km/h & 36 & 72 & 108 \\
m/s & 10 & 20 & 30 \\
\end{tblr}\par\vspace{2pt}
{\footnotesize\sffamily \R{上}{うえ}の\R{段}{だん}÷3.6＝\R{下}{した}の\R{段}{だん}\qquad \R{下}{した}の\R{段}{だん}×3.6＝\R{上}{うえ}の\R{段}{だん}}
\end{bkcenter}
÷3.6を\R{忘}{わす}れたときも、\R{単位}{たん|い}つきの1をかければ、いつでも\R{作}{つく}り\R{直}{なお}せます。\R{秒速}{びょう|そく}の\R{数}{かず}は\R{時速}{じ|そく}の\R{数}{かず}より\R{小}{ちい}さくなるはずです。1\R{秒}{びょう}は1\R{時間}{じ|かん}よりずっと\R{短}{みじか}く、\R{進}{すす}む\R{道}{みち}のりも\R{短}{みじか}いからです（ただしmで\R{数}{かぞ}えるので、3600\R{分}{ぶん}の1ではなく3.6\R{分}{ぶん}の1になる）。

\bksubsection{8}{\R{面積}{めん|せき}・\R{体積}{たい|せき}の\R{換算}{かん|さん}}

\noindent\begin{minipage}[c]{0.6\linewidth}
\leavevmode\label{bk:2-10:menseki}\qty{1}{m^2}は、1\R{辺}{ぺん}\qty{1}{m}の\R{正方形}{せい|ほう|けい}の\R{面積}{めん|せき}です。\R{右}{みぎ}の\R{図}{ず}のように、1\R{辺}{ぺん}は\qty{100}{cm}なので、
\[ 1\,\mathrm{m^2}=100\,\mathrm{cm}\times100\,\mathrm{cm}=10000\,\mathrm{cm^2} \]
\R{長}{なが}さが100\R{倍}{ばい}になると、\R{面積}{めん|せき}は $100\times100$ で10000\R{倍}{ばい}です。\R{図}{ず}の1マス（\qty{10}{cm}×\qty{10}{cm}＝\qty{100}{cm^2}）が100こある、と\R{数}{かぞ}えても\R{同}{おな}じです。
\end{minipage}\hfill
\begin{minipage}[c]{0.37\linewidth}\centering
\begin{tikzpicture}[x=1mm, y=1mm]
  \fill[bkBlue!12] (0,0) rectangle (34,34);
  \foreach \k in {1,...,9} {\draw[bkBlue!40, line width=0.3pt] (\k*3.4,0) -- (\k*3.4,34) (0,\k*3.4) -- (34,\k*3.4);}
  \fill[bkOrange!70] (0,30.6) rectangle (3.4,34);
  \draw[bkBlue!70!black, line width=0.8pt] (0,0) rectangle (34,34);
  \node[anchor=north, font=\sffamily\footnotesize, inner sep=0pt] at (17,-1.8) {\qty{100}{cm}（\qty{1}{m}）};
  \node[anchor=south, font=\sffamily\footnotesize, inner sep=0pt, rotate=90] at (-1.8,17) {\qty{100}{cm}};
  \draw[bkOrange!80!black, line width=0.6pt] (3.6,32.3) -- (5.6,37.6);
  \node[anchor=west, font=\sffamily\footnotesize, inner sep=0pt] at (6.2,38.6) {1マス＝10\,cm×10\,cm};
\end{tikzpicture}
\end{minipage}

\vspace{3pt}
\R{単位}{たん|い}つきの1で\R{書}{か}くと、$\mathrm{m^2}$ は $\mathrm{m}\times\mathrm{m}$ でmが2つあるので、\qty{100}{cm/m}を2\R{回}{かい}かけます。
\[ 1\,\mathrm{m^2}=1\,\bkiikc{\mathrm{m}}\times\bkiikc{\mathrm{m}}\times\frac{100\,\mathrm{cm}}{1\,\bkiikc{\mathrm{m}}}\times\frac{100\,\mathrm{cm}}{1\,\bkiikc{\mathrm{m}}}=10000\,\mathrm{cm^2} \]
\R{体積}{たい|せき}の $\mathrm{m^3}$ はmが3つなので、\qty{100}{cm/m}を3\R{回}{かい}かけて \mbox{\qty{1}{m^3}＝\qty{1000000}{cm^3}}。\R{体積}{たい|せき}では、\mbox{\qty{1}{L}＝\qty{1000}{cm^3}}、\mbox{\qty{1}{mL}＝\qty{1}{cm^3}}（2-7）もよく\R{使}{つか}います。\qty{1000}{cm^3/L}や\qty{1000}{mL/L}も、\R{単位}{たん|い}つきの1です。

\bksubsection{9}{\R{単位}{たん|い}も\R{計算}{けい|さん}できる}

\leavevmode\label{bk:2-10:keisan}\R{自転車}{じ|てん|しゃ}が、\R{止}{と}まった\R{状態}{じょう|たい}からこぎ\R{出}{だ}して、4\R{秒}{びょう}で\R{秒速}{びょう|そく}\qty{6}{m}になりました。\R{速}{はや}さが\R{同}{おな}じ\R{分}{ぶん}ずつふえたとすると、1\R{秒}{びょう}ごとに\R{速}{はや}さは\R{何}{なん}m/sずつふえたかな？

\noindent\begin{minipage}[c]{0.5\linewidth}
\R{速}{はや}さの\R{変化量}{へん|か|りょう}（2-3）は、$6-0=6$ で\qty{6}{m/s}。これを4\R{秒}{びょう}に\R{等}{ひと}しく\R{分}{わ}けると、1\R{秒}{びょう}あたりのふえ\R{方}{かた}がわかります。\R{数}{かず}は\R{数}{かず}どうし、\R{単位}{たん|い}は\R{単位}{たん|い}どうしで\R{計算}{けい|さん}します。
\end{minipage}\hfill
\begin{minipage}[c]{0.47\linewidth}\centering
\begin{tikzpicture}[x=1mm, y=1mm]
  \draw[bkInk!70, line width=0.7pt] (0,0) -- (56,0);
  \foreach \t/\v in {0/0,1/1.5,2/3,3/4.5,4/6} {
    \fill[bkBlue!45] ({\t*12+1.5},0) rectangle ({\t*12+8.5},{\v*4.2});
    \node[anchor=north, font=\sffamily\footnotesize, inner sep=0pt] at ({\t*12+5},-1.6) {\t\R{秒}{びょう}};
    \node[anchor=south, font=\sffamily\footnotesize, inner sep=0pt] at ({\t*12+5},{\v*4.2+1.4}) {\v};}
  \foreach \t/\v in {0/0,1/1.5,2/3,3/4.5} {\node[font=\sffamily\bfseries\footnotesize, text=bkRed, inner sep=0pt] at ({\t*12+11},{\v*4.2+3.2}) {＋};}
  \node[anchor=west, font=\sffamily\footnotesize, text=bkGray, inner sep=0pt] at (-1,33) {\R{速}{はや}さ（m/s）};
\end{tikzpicture}
\bkfigcaption{1\R{秒}{びょう}ごとに\qty{1.5}{m/s}ずつ\R{速}{はや}くなる。}
\end{minipage}
\[ \frac{6\,\mathrm{m/s}}{4\,\mathrm{s}}=\frac{6}{4}\times\frac{\mathrm{m/s}}{\mathrm{s}}=1.5\,\mathrm{m/s^2} \]
\R{単位}{たん|い}は $(\mathrm{m/s})\div\mathrm{s}$ です。m/sは $\dfrac{\mathrm{m}}{\mathrm{s}}$ なので、さらにsでわると\R{分母}{ぶん|ぼ}にsが2つならびます。
\[ (\mathrm{m/s})\div\mathrm{s}=\frac{\mathrm{m}}{\mathrm{s}\times\mathrm{s}}=\frac{\mathrm{m}}{\mathrm{s^2}}=\mathrm{m/s^2} \]
$\mathrm{s}\times\mathrm{s}$ を $\mathrm{s^2}$ と\R{書}{か}くのは、2-7の\R{累乗}{るい|じょう}と\R{同}{おな}じです。
\unit{m/s^2}は「メートル\R{毎秒}{まい|びょう}\R{毎秒}{まい|びょう}」と\R{読}{よ}み、1\R{秒}{びょう}ごとに\R{速}{はや}さが\R{何}{なん}m/sふえるかを\R{表}{あらわ}します。\R{上}{うえ}のぼうグラフでも、1\R{秒}{びょう}ごとに1.5ずつ\R{高}{たか}くなっています。\R{落}{お}ちるボールは、1\R{秒}{びょう}ごとに\R{約}{やく}\qty{10}{m/s}ずつ\R{速}{はや}くなりました（2-0）。これは\R{約}{やく}\qty{10}{m/s^2}です（もっとくわしくは\qty{9.8}{m/s^2}）。

\bksubsection{10}{\R{単位}{たん|い}のチェック}

\leavevmode\label{bk:2-10:check}\R{式}{しき}の\R{両側}{りょう|がわ}の\R{単位}{たん|い}は、いつも\R{同}{おな}じになります。\R{左}{ひだり}と\R{右}{みぎ}は\R{同}{おな}じ\R{量}{りょう}を\R{表}{あらわ}しているので、\R{量}{りょう}の\R{種類}{しゅ|るい}も\R{同}{おな}じはずだからです。「\R{道}{みち}のり＝\R{速}{はや}さ×\R{時間}{じ|かん}」で\R{確}{たし}かめましょう。
\[ \frac{\mathrm{m}}{\bkiikc{\mathrm{s}}}\times\bkiikc{\mathrm{s}}=\mathrm{m} \]
\R{右側}{みぎ|がわ}の\R{単位}{たん|い}がmになり、\R{左側}{ひだり|がわ}の\R{道}{みち}のりのmと\R{合}{あ}いました。
もし「\R{道}{みち}のり＝\R{速}{はや}さ÷\R{時間}{じ|かん}」とまちがえて\R{覚}{おぼ}えていたら、\R{右}{みぎ}は $(\mathrm{m/s})\div\mathrm{s}=\mathrm{m/s^2}$ になり、\R{左}{ひだり}のmと\R{合}{あ}いません。\R{単位}{たん|い}が\R{合}{あ}わない\R{式}{しき}は、どこかがまちがっています。このように\R{単位}{たん|い}で\R{式}{しき}を\R{確}{たし}かめることを、\textbf{\R{単位}{たん|い}のチェック}といいます。\bkidx{たんいのちぇっく}{\R{単位}{たん|い}のチェック}\R{単位}{たん|い}が\R{合}{あ}っても、\R{数}{かず}がまちがっていることはあるので、\R{単位}{たん|い}のチェックは「まちがいを\R{見}{み}つける\R{道具}{どう|ぐ}」です。

\begin{taisetsu}
\qty{1}{km}＝\qty{1000}{m}、\qty{1}{h}＝\qty{3600}{s}。\R{時速}{じ|そく}を\R{秒速}{びょう|そく}にするには÷3.6\\
1\R{分}{ぷん}＝60\R{秒}{びょう}、1\R{時間}{じ|かん}＝3600\R{秒}{びょう}。\R{単位}{たん|い}をそろえてから\R{計算}{けい|さん}する。\\
\R{換算}{かん|さん}は「\R{単位}{たん|い}つきの1」をかける：\qty{2}{min}×\qty{60}{s/min}＝\qty{120}{s}（minが\R{約分}{やく|ぶん}されて\R{消}{き}える）\\
\R{単位}{たん|い}も\R{計算}{けい|さん}できる。(\unit{m/s})÷\unit{s}＝\unit{m/s^2}。\R{式}{しき}の\R{両側}{りょう|がわ}の\R{単位}{たん|い}は\R{同}{おな}じ。
\end{taisetsu}

\begin{kotoba}
\small
\bkword{SI・\R{基本}{き|ほん}\R{単位}{たん|い}}{\R{世界}{せ|かい}\R{共通}{きょう|つう}の\R{単位}{たん|い}のきまり。もとになる7つの\R{単位}{たん|い}。}
\bkword{\R{組}{く}み\R{立}{た}てた\R{単位}{たん|い}}{\R{基本}{き|ほん}\R{単位}{たん|い}をかけたりわったりした\R{単位}{たん|い}（\unit{m^2}、\unit{m/s}）。}
\bkword{\R{接頭語}{せっ|とう|ご}}{\R{単位}{たん|い}の\R{前}{まえ}につけて、\R{大}{おお}きさを\R{表}{あらわ}す\R{記号}{き|ごう}（k、m、{\textmu}）。}
\bkword{\R{単位}{たん|い}つきの1}{\R{上}{うえ}と\R{下}{した}が\R{同}{おな}じ\R{大}{おお}きさの\R{量}{りょう}の\R{分数}{ぶん|すう}。\qty{60}{s/min}、\qty{1000}{m/km}など。}
\bkword{\R{換算}{かん|さん}}{\R{単位}{たん|い}つきの1をかけて、\R{量}{りょう}を\R{変}{か}えずに\R{単位}{たん|い}だけを\R{取}{と}りかえること。}
\end{kotoba}

%% ---------------- 例題 ----------------
\bksection{reidai}{\R{例題}{れい|だい}でたしかめよう}

\begin{reidai}[\R{接頭語}{せっ|とう|ご}のついた\R{単位}{たん|い}を\R{換算}{かん|さん}する]
（1）\qty{2.5}{km}は\R{何}{なん}m？\quad（2）\qty{350}{mg}は\R{何}{なん}g？\par
\kangaekata \qty{1000}{m/km}と\qty{1000}{mg/g}（1\,gあたり1000\,mg）を、\R{消}{け}したい\R{単位}{たん|い}が\R{分母}{ぶん|ぼ}にくる\R{向}{む}きでかけます。\par
\shiki（1）kmを\R{消}{け}す。$2.5\,\bkiikc{\mathrm{km}}\times\dfrac{1000\,\mathrm{m}}{1\,\bkiikc{\mathrm{km}}}=2.5\times1000\,\mathrm{m}=2500\,\mathrm{m}$\par
\hspace*{2.6\zw}（2）mgを\R{消}{け}す。$350\,\bkiikc{\mathrm{mg}}\times\dfrac{1\,\mathrm{g}}{1000\,\bkiikc{\mathrm{mg}}}=350\div1000\,\mathrm{g}=0.35\,\mathrm{g}$\par
\kotae（1）\qty{2500}{m}\quad（2）\qty{0.35}{g}
\end{reidai}

\begin{jibun}
\R{単位}{たん|い}つきの1をかけて\R{換算}{かん|さん}しましょう。（1）\qty{0.8}{km}は\R{何}{なん}m？\quad（2）\qty{1.2}{kg}は\R{何}{なん}g？\quad（3）\qty{2500}{mg}は\R{何}{なん}g？
\bkjans{（1）\nobreak \qty{800}{m}\quad（2）\nobreak \qty{1200}{g}\quad（3）\nobreak \qty{2.5}{g}}
\end{jibun}

\begin{reidai}[\R{時間}{じ|かん}の\R{単位}{たん|い}をつないで\R{換算}{かん|さん}する]
（1）1\R{日}{にち}は\R{何}{なん}\R{秒}{びょう}ですか。\quad（2）1\R{年}{ねん}（365\R{日}{にち}）は\R{何}{なん}\R{時間}{じ|かん}ですか。\par
\kangaekata \pageref{bk:2-10:jikan}ページの\R{図}{ず}の\R{順}{じゅん}に、\R{単位}{たん|い}つきの1を1つずつかけます。\par
\shiki（1）$1\,\bkiikc{\mathrm{d}}\times\dfrac{24\,\bkiikc{\mathrm{h}}}{1\,\bkiikc{\mathrm{d}}}\times\dfrac{60\,\bkiikc{\mathrm{min}}}{1\,\bkiikc{\mathrm{h}}}\times\dfrac{60\,\mathrm{s}}{1\,\bkiikc{\mathrm{min}}}$\par
\hspace*{4.6\zw}$=24\times60\times60\,\mathrm{s}=1440\times60\,\mathrm{s}=86400\,\mathrm{s}$\par
\hspace*{2.6\zw}（2）1\,\bkiikt{\R{年}{ねん}}×365\,\R{日}{にち}/\bkiikt{\R{年}{ねん}}＝365\,\R{日}{にち}。365\R{日}{にち}は\qty{365}{d}と\R{書}{か}ける。\par
\hspace*{4.6\zw}$365\,\bkiikc{\mathrm{d}}\times\dfrac{24\,\mathrm{h}}{1\,\bkiikc{\mathrm{d}}}=365\times24\,\mathrm{h}=8760\,\mathrm{h}$\par
\kotae（1）86400\R{秒}{びょう}\quad（2）8760\R{時間}{じ|かん}
\end{reidai}

\begin{jibun}
（1）2\R{時間}{じ|かん}は\R{何}{なん}\R{秒}{びょう}ですか。\quad（2）1\R{週間}{しゅう|かん}（7\R{日}{にち}）は\R{何}{なん}\R{分}{ぷん}ですか。どちらも\R{単位}{たん|い}つきの1をかける\R{式}{しき}で\R{書}{か}きましょう。
\bkjans{（1）\nobreak \qty{2}{h}×\qty{60}{min/h}×\qty{60}{s/min}＝7200\R{秒}{びょう}\quad（2）\nobreak \qty{7}{d}×\qty{24}{h/d}×\qty{60}{min/h}＝10080\R{分}{ぷん}}
\end{jibun}

\begin{reidai}[\R{時間}{じ|かん}の\R{単位}{たん|い}をそろえてから、\R{速}{はや}さを\R{求}{もと}める]
\R{自転車}{じ|てん|しゃ}で、1\R{分}{ぷん}30\R{秒}{びょう}に\qty{450}{m}\R{進}{すす}みました。\R{速}{はや}さは\R{何}{なん}m/sですか。\par
\kangaekata m/sで\R{答}{こた}えるので、\R{時間}{じ|かん}を\R{先}{さき}に\R{秒}{びょう}にそろえます。1\R{分}{ぷん}を\R{秒}{びょう}に\R{換算}{かん|さん}してから、30\R{秒}{びょう}をたします。\par
\shiki $1\,\bkiikc{\mathrm{min}}\times\dfrac{60\,\mathrm{s}}{1\,\bkiikc{\mathrm{min}}}=60\,\mathrm{s}$。$60\,\mathrm{s}+30\,\mathrm{s}=90\,\mathrm{s}$\par
\hspace*{2.6\zw}\R{速}{はや}さ＝\R{道}{みち}のり÷\R{時間}{じ|かん}（2-3）。$450\,\mathrm{m}\div90\,\mathrm{s}=5\,\mathrm{m/s}$\par
\kotae \qty{5}{m/s}
\end{reidai}

\begin{jibun}
\R{走}{はし}って、2\R{分}{ふん}で\qty{600}{m}\R{進}{すす}みました。\R{速}{はや}さは\R{何}{なん}m/sですか。
\bkjans{\qty{2}{min}×\qty{60}{s/min}＝\qty{120}{s}。$600\div120=5$ で\qty{5}{m/s}}
\end{jibun}

\begin{reidai}[\R{時速}{じ|そく}と\R{秒速}{びょう|そく}を\R{換算}{かん|さん}する]
（1）\R{時速}{じ|そく}\qty{54}{km}は\R{秒速}{びょう|そく}\R{何}{なん}m？\quad（2）\R{秒速}{びょう|そく}\qty{25}{m}は\R{時速}{じ|そく}\R{何}{なん}km？\par
\kangaekata kmとmは\qty{1000}{m/km}で、hとsは\qty{3600}{s/h}でつなぎます。\R{消}{け}したい\R{単位}{たん|い}が\R{約分}{やく|ぶん}される\R{向}{む}きにかけます。\R{最後}{さい|ご}に÷3.6・×3.6で\R{確}{たし}かめます。\par
\shiki（1）$54\,\dfrac{\bkiikc{\mathrm{km}}}{\bkiikc{\mathrm{h}}}\times\dfrac{1000\,\mathrm{m}}{1\,\bkiikc{\mathrm{km}}}\times\dfrac{1\,\bkiikc{\mathrm{h}}}{3600\,\mathrm{s}}=\dfrac{54000}{3600}\,\mathrm{m/s}=15\,\mathrm{m/s}$\par
\hspace*{2.6\zw}（2）$25\,\dfrac{\bkiikc{\mathrm{m}}}{\bkiikc{\mathrm{s}}}\times\dfrac{3600\,\bkiikc{\mathrm{s}}}{1\,\mathrm{h}}\times\dfrac{1\,\mathrm{km}}{1000\,\bkiikc{\mathrm{m}}}=\dfrac{90000}{1000}\,\mathrm{km/h}=90\,\mathrm{km/h}$\par
\hspace*{2.6\zw}\R{確}{たし}かめ：$54\div3.6=15$、\ $25\times3.6=90$\par
\kotae（1）\qty{15}{m/s}\quad（2）\qty{90}{km/h}
\end{reidai}

\begin{jibun}
（1）\R{時速}{じ|そく}\qty{18}{km}は\R{秒速}{びょう|そく}\R{何}{なん}m？\quad（2）\R{秒速}{びょう|そく}\qty{30}{m}は\R{時速}{じ|そく}\R{何}{なん}km？
\bkjans{（1）\nobreak \qty{5}{m/s}\quad（2）\nobreak \qty{108}{km/h}}
\end{jibun}

\begin{motto}[2]
\textbf{\R{早見表}{はや|み|ひょう}：よく\R{使}{つか}う「\R{単位}{たん|い}つきの1」}\par
どれも「1\R{右}{みぎ}の\R{単位}{たん|い}あたり、\R{左}{ひだり}の\R{単位}{たん|い}がいくつか」を\R{表}{あらわ}します。\R{上下}{じょう|げ}を\R{入}{い}れかえても1なので、\R{消}{け}したい\R{単位}{たん|い}が\R{分母}{ぶん|ぼ}にくる\R{向}{む}きで\R{使}{つか}います。
\begin{bkcenter}
\begin{bktable}{colspec={Q[l,m,wd=21mm] Q[l,m,wd=30mm] Q[l,m,wd=20mm] Q[l,m,wd=37mm]}, rowsep=1pt, colsep=4pt, cells={font=\small}, row{1}={font=\sffamily\bfseries\small, bg=bkGray!12}}
\R{長}{なが}さ・\R{重}{おも}さ & \R{意味}{い|み} & \R{時間}{じ|かん}・\R{体積}{たい|せき} & \R{意味}{い|み} \\
\qty{1000}{m/km} & 1\,kmあたり1000\,m & \qty{60}{s/min} & 1\R{分}{ぷん}あたり60\R{秒}{びょう} \\
\qty{100}{cm/m} & 1\,mあたり100\,cm & \qty{60}{min/h} & 1\R{時間}{じ|かん}あたり60\R{分}{ぷん} \\
\qty{10}{mm/cm} & 1\,cmあたり10\,mm & \qty{3600}{s/h} & 1\R{時間}{じ|かん}あたり3600\R{秒}{びょう} \\
\qty{1000}{g/kg} & 1\,kgあたり1000\,g & \qty{24}{h/d} & 1\R{日}{にち}あたり24\R{時間}{じ|かん} \\
\qty{1000}{mg/g} & 1\,gあたり1000\,mg & \qty{1000}{mL/L} & 1\,Lあたり1000\,mL \\
\end{bktable}
\end{bkcenter}
\end{motto}

\begin{reidai}[\R{速}{はや}さと\R{時間}{じ|かん}の\R{単位}{たん|い}がちがうとき、\R{道}{みち}のりを\R{求}{もと}める]
（1）\R{時速}{じ|そく}\qty{54}{km}で20\R{分}{ぷん}\R{走}{はし}ると、\R{何}{なん}km\R{進}{すす}みますか。\quad（2）\R{秒速}{びょう|そく}\qty{15}{m}で1\R{時間}{じ|かん}\R{走}{はし}ると、\R{何}{なん}km\R{進}{すす}みますか。\par
\kangaekata \R{道}{みち}のり＝\R{速}{はや}さ×\R{時間}{じ|かん}（2-3）。\R{速}{はや}さの\R{分母}{ぶん|ぼ}の\R{単位}{たん|い}と、\R{時間}{じ|かん}の\R{単位}{たん|い}がちがうので、\R{単位}{たん|い}つきの1をはさんで\R{約分}{やく|ぶん}します。\par
\shiki（1）$54\,\dfrac{\mathrm{km}}{\bkiikc{\mathrm{h}}}\times20\,\bkiikc{\mathrm{min}}\times\dfrac{1\,\bkiikc{\mathrm{h}}}{60\,\bkiikc{\mathrm{min}}}=\dfrac{54\times20}{60}\,\mathrm{km}=18\,\mathrm{km}$\par
\hspace*{2.6\zw}（2）$15\,\dfrac{\mathrm{m}}{\bkiikc{\mathrm{s}}}\times1\,\bkiikc{\mathrm{h}}\times\dfrac{3600\,\bkiikc{\mathrm{s}}}{1\,\bkiikc{\mathrm{h}}}=54000\,\mathrm{m}$\par
\hspace*{4.6\zw}$54000\,\bkiikc{\mathrm{m}}\times\dfrac{1\,\mathrm{km}}{1000\,\bkiikc{\mathrm{m}}}=54\,\mathrm{km}$\par
\kotae（1）\qty{18}{km}\quad（2）\qty{54}{km}
\end{reidai}

\begin{jibun}
（1）\R{時速}{じ|そく}\qty{60}{km}で45\R{分}{ふん}\R{走}{はし}ると、\R{何}{なん}km\R{進}{すす}みますか。\quad（2）\R{秒速}{びょう|そく}\qty{5}{m}で10\R{分}{ぷん}\R{走}{はし}ると、\R{何}{なん}km\R{進}{すす}みますか。
\bkjans{（1）\nobreak \qty{60}{km/h}×\qty{45}{min}×1\,h/60\,min＝\qty{45}{km}\quad（2）\nobreak \qty{5}{m/s}×\qty{600}{s}＝\qty{3000}{m}＝\qty{3}{km}}
\end{jibun}

\begin{reidai}[\R{単位}{たん|い}をそろえて、\R{何}{なん}こ\R{分}{ぶん}・\R{何倍}{なん|ばい}かを\R{求}{もと}める]
（1）\qty{1.5}{L}のお\R{茶}{ちゃ}を、\qty{250}{mL}ずつコップに\R{分}{わ}けます。\R{何}{なん}ばいに\R{分}{わ}けられますか。\par
（2）\R{時速}{じ|そく}\qty{36}{km}で\R{走}{はし}る\R{自転車}{じ|てん|しゃ}は、\R{秒速}{びょう|そく}\qty{5}{m}で\R{走}{はし}る\R{人}{ひと}の\R{何倍}{なん|ばい}の\R{速}{はや}さですか。\par
\kangaekata \R{単位}{たん|い}をそろえてからわります。\R{同}{おな}じ\R{単位}{たん|い}どうしでわると、\R{単位}{たん|い}は\R{約分}{やく|ぶん}されて\R{消}{き}え、「\R{何}{なん}こ\R{分}{ぶん}」「\R{何倍}{なん|ばい}」を\R{表}{あらわ}す\R{数}{かず}だけが\R{残}{のこ}ります。\par
\shiki（1）$1.5\,\bkiikc{\mathrm{L}}\times\dfrac{1000\,\mathrm{mL}}{1\,\bkiikc{\mathrm{L}}}=1500\,\mathrm{mL}$。$\dfrac{1500\,\bkiikc{\mathrm{mL}}}{250\,\bkiikc{\mathrm{mL}}}=6$\par
\hspace*{2.6\zw}（2）$36\div3.6=10$ で\qty{10}{m/s}。$\dfrac{10\,\bkiikc{\mathrm{m/s}}}{5\,\bkiikc{\mathrm{m/s}}}=2$\par
\kotae（1）6ぱい\quad（2）2\R{倍}{ばい}
\end{reidai}

\begin{jibun}
（1）\qty{3}{L}の\R{水}{みず}を、\qty{150}{mL}ずつ\R{分}{わ}けます。\R{何}{なん}ばいに\R{分}{わ}けられますか。\quad（2）\R{時速}{じ|そく}\qty{54}{km}の\R{車}{くるま}は、\R{秒速}{びょう|そく}\qty{5}{m}の\R{自転車}{じ|てん|しゃ}の\R{何倍}{なん|ばい}の\R{速}{はや}さですか。
\bkjans{（1）\nobreak 3000\,mL÷150\,mL＝20で20ぱい\quad（2）\nobreak \qty{15}{m/s}÷\qty{5}{m/s}＝3で3\R{倍}{ばい}}
\end{jibun}

\begin{jibun}
（1）\qty{2}{L}の\R{牛乳}{ぎゅう|にゅう}を、\qty{400}{mL}ずつ\R{分}{わ}けます。\R{何}{なん}ばいに\R{分}{わ}けられますか。\quad（2）\R{秒速}{びょう|そく}\qty{25}{m}の\R{電車}{でん|しゃ}は、\R{時速}{じ|そく}\qty{18}{km}の\R{自転車}{じ|てん|しゃ}の\R{何倍}{なん|ばい}の\R{速}{はや}さですか。
\bkjans{（1）\nobreak 2000\,mL÷400\,mL＝5で5はい\quad（2）\nobreak \qty{18}{km/h}は\qty{5}{m/s}。$25\div5=5$ で5\R{倍}{ばい}}
\end{jibun}

\begin{reidai}[\R{単位}{たん|い}を\R{計算}{けい|さん}する]
（1）\R{止}{と}まっていた\R{電車}{でん|しゃ}が、20\R{秒}{びょう}で\qty{15}{m/s}になりました。1\R{秒}{びょう}ごとに\R{速}{はや}さはいくつふえましたか（\R{同}{おな}じ\R{分}{ぶん}ずつふえたとする）。\par
（2）「\R{道}{みち}のり＝\R{速}{はや}さ÷\R{時間}{じ|かん}」は\R{正}{ただ}しいですか。\R{単位}{たん|い}のチェックをしましょう。\par
\kangaekata（1）\R{速}{はや}さの\R{変化量}{へん|か|りょう}÷\R{時間}{じ|かん}。\R{数}{かず}と\R{単位}{たん|い}を\R{別}{べつ}べつに\R{計算}{けい|さん}します。（2）\R{右}{みぎ}の\R{単位}{たん|い}がmになるか\R{見}{み}ます。\par
\shiki（1）\R{変化量}{へん|か|りょう}は $15-0=15$ で\qty{15}{m/s}。$15\,\mathrm{m/s}\div20\,\mathrm{s}=0.75\,\mathrm{m/s^2}$\par
\hspace*{2.6\zw}（2）$(\mathrm{m/s})\div\mathrm{s}=\mathrm{m/s^2}$。\R{道}{みち}のりの\R{単位}{たん|い}mと\R{合}{あ}わない\par
\kotae（1）\qty{0.75}{m/s^2}\quad（2）まちがい（\R{正}{ただ}しくは\R{速}{はや}さ×\R{時間}{じ|かん}）
\end{reidai}

\begin{jibun}
（1）\R{止}{と}まった\R{車}{くるま}が5\R{秒}{びょう}で\qty{10}{m/s}になりました。1\R{秒}{びょう}ごとに\R{速}{はや}さはいくつふえましたか。（2）「\R{時間}{じ|かん}＝\R{道}{みち}のり÷\R{速}{はや}さ」の\R{単位}{たん|い}を\R{調}{しら}べましょう。
\bkjans{（1）\nobreak \qty{2}{m/s^2}\quad（2）\nobreak m÷(m/s)＝sで、\R{時間}{じ|かん}の\R{単位}{たん|い}と\R{合}{あ}う（\R{正}{ただ}しい）}
\end{jibun}

\begin{jibun}
（1）\R{止}{と}まっていた\R{自転車}{じ|てん|しゃ}が、4\R{秒}{びょう}で\qty{8}{m/s}になりました。1\R{秒}{びょう}ごとに\R{速}{はや}さはいくつふえましたか。（2）「\R{速}{はや}さ＝\R{道}{みち}のり×\R{時間}{じ|かん}」は\R{正}{ただ}しいですか。\R{単位}{たん|い}のチェックをしましょう。
\bkjans{（1）\nobreak \qty{2}{m/s^2}\quad（2）\nobreak m×s で、\R{速}{はや}さの\R{単位}{たん|い}m/sと\R{合}{あ}わない（まちがい）}
\end{jibun}

\begin{reidai}[\R{面積}{めん|せき}・\R{体積}{たい|せき}・{\textmu}mの\R{換算}{かん|さん}]
（1）\qty{2}{m^2}は\R{何}{なん}\unit{cm^2}？\quad（2）\qty{2.5}{L}は\R{何}{なん}\unit{cm^3}？\quad（3）\qty{70}{\micro m}は\R{何}{なん}mm？\quad（4）\qty{0.3}{mm}は\R{何}{なん}\unit{\micro m}？\par
\kangaekata \qty{100}{cm/m}（\unit{m^2}では2\R{回}{かい}）、\qty{1000}{cm^3/L}、\qty{1000}{\micro m/mm}を\R{使}{つか}います。\par
\shiki（1）$2\,\mathrm{m^2}\times\dfrac{100\,\mathrm{cm}}{1\,\mathrm{m}}\times\dfrac{100\,\mathrm{cm}}{1\,\mathrm{m}}=2\times10000\,\mathrm{cm^2}=20000\,\mathrm{cm^2}$\par
\hspace*{2.6\zw}（2）$2.5\,\bkiikc{\mathrm{L}}\times\dfrac{1000\,\mathrm{cm^3}}{1\,\bkiikc{\mathrm{L}}}=2500\,\mathrm{cm^3}$\par
\hspace*{2.6\zw}（3）$70\,\bkiikc{\unit{\micro m}}\times\dfrac{1\,\mathrm{mm}}{1000\,\bkiikc{\unit{\micro m}}}=0.07\,\mathrm{mm}$\par
\hspace*{2.6\zw}（4）$0.3\,\bkiikc{\mathrm{mm}}\times\dfrac{1000\,\unit{\micro m}}{1\,\bkiikc{\mathrm{mm}}}=300\,\unit{\micro m}$\par
\kotae（1）\qty{20000}{cm^2}\quad（2）\qty{2500}{cm^3}\quad（3）\qty{0.07}{mm}\quad（4）\qty{300}{\micro m}
\end{reidai}

\begin{jibun}
（1）\qty{3}{m^2}は\R{何}{なん}\unit{cm^2}？\quad（2）\qty{0.3}{L}は\R{何}{なん}\unit{cm^3}？\quad（3）\qty{500}{\micro m}は\R{何}{なん}mm？\quad（4）\qty{0.04}{mm}は\R{何}{なん}\unit{\micro m}？
\bkjans{（1）\nobreak \qty{30000}{cm^2}\quad（2）\nobreak \qty{300}{cm^3}\quad（3）\nobreak \qty{0.5}{mm}\quad（4）\nobreak \qty{40}{\micro m}}
\end{jibun}

\begin{jibun}
（1）\qty{0.5}{m^2}は\R{何}{なん}\unit{cm^2}？\quad（2）\qty{1.2}{L}は\R{何}{なん}\unit{cm^3}？\quad（3）\qty{2}{mm}は\R{何}{なん}\unit{\micro m}？\quad（4）\qty{250}{\micro m}は\R{何}{なん}mm？
\bkjans{（1）\nobreak \qty{5000}{cm^2}\quad（2）\nobreak \qty{1200}{cm^3}\quad（3）\nobreak \qty{2000}{\micro m}\quad（4）\nobreak \qty{0.25}{mm}}
\end{jibun}

\clearpage
%% ---------------- まちがえやすいところ ----------------
\bksection{tsumazuki}{まちがえやすいところをチェック}

\begin{tsumazuki}{\R{時速}{じ|そく}72\,kmを\R{秒速}{びょう|そく}にするとき、3.6をかけてしまう}
\bkng{$72\times3.6=259.2$ で\qty{259.2}{m/s}}
\bkok{$72\div3.6=20$ で\qty{20}{m/s}}
1\R{秒}{びょう}は1\R{時間}{じ|かん}よりずっと\R{短}{みじか}いので、\R{秒速}{びょう|そく}の\R{数}{かず}は\R{時速}{じ|そく}より\R{小}{ちい}さくなるはずです。\qty{259.2}{m/s}は、\R{飛行機}{ひ|こう|き}くらいの\R{速}{はや}さです。
\end{tsumazuki}

\begin{tsumazuki}{\R{単位}{たん|い}つきの1を、\R{上下}{じょう|げ}さかさまにかけてしまう}
\bkng{$3\,\mathrm{min}\times\dfrac{1\,\mathrm{min}}{60\,\mathrm{s}}=\dfrac{3}{60}\,\dfrac{\mathrm{min}\times\mathrm{min}}{\mathrm{s}}$}
\bkok{$3\,\bkiikc{\mathrm{min}}\times\dfrac{60\,\mathrm{s}}{1\,\bkiikc{\mathrm{min}}}=180\,\mathrm{s}$}
\R{消}{け}したいminが\R{分母}{ぶん|ぼ}にくる\R{向}{む}きでかけます。さかさまだと、minが\R{消}{き}えずに\R{単位}{たん|い}が $\dfrac{\mathrm{min}\times\mathrm{min}}{\mathrm{s}}$ になり、\R{時間}{じ|かん}の\R{単位}{たん|い}になりません。
\end{tsumazuki}

\begin{tsumazuki}{\R{面積}{めん|せき}の\R{換算}{かん|さん}で100\R{倍}{ばい}にする}
\bkng{\qty{1}{m^2}＝\qty{100}{cm^2}\quad（\qty{1}{m}＝\qty{100}{cm}だから？）}
\bkok{\qty{1}{m^2}＝$100\,\mathrm{cm}\times100\,\mathrm{cm}$＝\qty{10000}{cm^2}}
\R{面積}{めん|せき}は\R{長}{なが}さを2\R{回}{かい}かけた\R{量}{りょう}なので、\qty{100}{cm/m}も2\R{回}{かい}かけます。
\end{tsumazuki}

\begin{tsumazuki}{\R{秒速}{びょう|そく}5\,mで2\R{分}{ふん}。\R{単位}{たん|い}をそろえずに\R{計算}{けい|さん}する}
\bkng{$5\times2=10$ で\qty{10}{m}}
\bkok{\qty{2}{min}×\qty{60}{s/min}＝\qty{120}{s}。$5\times120=600$ で\qty{600}{m}}
\unit{m/s}×minでは、sとminが\R{約分}{やく|ぶん}できず、mが\R{残}{のこ}りません。\R{単位}{たん|い}をそろえてから\R{計算}{けい|さん}します。
\end{tsumazuki}

\begin{tsumazuki}{1\R{時間}{じ|かん}を100\R{分}{ぷん}、1\R{分}{ぷん}を100\R{秒}{びょう}と\R{考}{かんが}えてしまう}
\bkng{1\R{時間}{じ|かん}30\R{分}{ぷん}＝1.3\R{時間}{じ|かん}}
\bkok{\qty{30}{min}÷\qty{60}{min/h}＝\qty{0.5}{h}。だから1\R{時間}{じ|かん}30\R{分}{ぷん}＝1.5\R{時間}{じ|かん}}
\R{長}{なが}さや\R{重}{おも}さは10・100・1000ずつ\R{単位}{たん|い}が\R{変}{か}わりますが、\R{時間}{じ|かん}は60ずつです。\R{時間}{じ|かん}は、\R{小数}{しょう|すう}の\R{位}{くらい}をそのまま\R{使}{つか}えません。
\end{tsumazuki}

\begin{column}{なぜ1\R{分}{ぷん}は60\R{秒}{びょう}なの？}
\R{長}{なが}さや\R{重}{おも}さは10\R{倍}{ばい}・1000\R{倍}{ばい}で\R{単位}{たん|い}が\R{変}{か}わるのに、\R{時間}{じ|かん}は60ずつです。これは、4000\R{年}{ねん}\R{以上}{い|じょう}\R{前}{まえ}の\R{古代}{こ|だい}メソポタミア（\R{今}{いま}のイラクのあたり）で、60ごとに\R{位}{くらい}が\R{上}{あ}がる\R{数}{かぞ}え\R{方}{かた}が\R{使}{つか}われていたなごりだといわれています。60は1・2・3・4・5・6・10・12・15・20・30・60でわり\R{切}{き}れるので、\R{半分}{はん|ぶん}（30\R{分}{ぷん}）、3\R{分}{ぶん}の1（20\R{分}{ぷん}）、4\R{分}{ぶん}の1（15\R{分}{ぷん}）がどれも\R{整数}{せい|すう}になり、\R{分}{わ}けやすいのです。\R{角度}{かく|ど}の1\R{周}{しゅう}360°も、60のなかまです。
\end{column}

%% ---------------- 物理・工学ではこう使う ----------------
\begin{butsuri}
\noindent\begin{minipage}[c]{0.6\linewidth}
\par\addvspace{3pt}\textbf{\R{単位}{たん|い}のとりちがえで、\R{火星}{か|せい}\R{探査機}{たん|さ|き}が\R{消}{き}えた}\par\noindent 1999\R{年}{ねん}、アメリカの\R{火星}{か|せい}\R{探査機}{たん|さ|き}「マーズ・クライメイト・オービター」が、\R{火星}{か|せい}に\R{着}{つ}いたところで\R{消}{き}えてしまいました。エンジンをふかす\R{強}{つよ}さを、あるチームはポンドという\R{単位}{たん|い}で、\R{別}{べつ}のチームはN（ニュートン）で\R{計算}{けい|さん}していたのです。1ポンドの\R{力}{ちから}は\R{約}{やく}4.45\,N。
\end{minipage}\hfill
\begin{minipage}[c]{0.37\linewidth}\centering
\begin{tikzpicture}[x=1mm, y=1mm]
  % 火星の地面の近く（高さはたてに 1 mm＝10 km）
  \fill[bkRed!55!bkOrange] (-22,-3) .. controls (-8,0.6) and (8,0.6) .. (22,-3) -- (22,-7) -- (-22,-7) -- cycle;
  \fill[bkOrange!25] (-22,-3) .. controls (-8,0.6) and (8,0.6) .. (22,-3) -- (22,5) .. controls (8,8.6) and (-8,8.6) .. (-22,5) -- cycle;
  \node[font=\sffamily\bfseries\footnotesize, text=white, inner sep=0pt] at (0,-4.4) {\R{火星}{か|せい}};
  \draw[bkGreen!60!black, line width=1pt, dash pattern=on 2pt off 1pt] (-22,12.6) .. controls (-8,16.2) and (8,16.2) .. (22,12.6);
  \draw[bkRed, line width=1.2pt] (-22,3.3) .. controls (-8,6.9) and (8,6.9) .. (22,3.3);
  \node[anchor=south, font=\sffamily\footnotesize, text=bkGreen!50!black, inner sep=0pt] at (0,16.6) {\R{予定}{よ|てい}　\R{高}{たか}さ\R{約}{やく}150\,km};
  \node[anchor=south, font=\sffamily\footnotesize, text=bkRed, fill=bkPurpleBg, inner sep=0.6pt] at (0,7.4) {\R{実際}{じっ|さい}　\R{約}{やく}57\,km};
\end{tikzpicture}
\end{minipage}
\par\vspace{2pt}
\R{計算}{けい|さん}がずれて、\R{探査機}{たん|さ|き}は\R{予定}{よ|てい}（\R{高}{たか}さ\R{約}{やく}\qty{150}{km}）よりずっと\R{低}{ひく}い、\R{約}{やく}\qty{57}{km}まで\R{火星}{か|せい}に\R{近}{ちか}づき、こわれてしまいました。\R{単位}{たん|い}をそろえて、\R{単位}{たん|い}のチェックをしていれば\R{防}{ふせ}げた\R{事故}{じ|こ}です。
\begin{itemize}[leftmargin=1.2\zw, itemsep=1pt, topsep=4pt, label={\color{bkPurple}●}]
\item \textbf{\R{薬}{くすり}のmgと\unit{\micro g}}\par\noindent \qty{1}{mg}は\qty{1000}{\micro g}です。\R{薬}{くすり}の\R{量}{りょう}の\R{単位}{たん|い}を\R{読}{よ}みまちがえると、1000\R{倍}{ばい}や1000\R{分}{ぶん}の1の\R{量}{りょう}になってしまいます。\R{薬}{くすり}をあつかう\R{人}{ひと}は、\R{数}{かず}といっしょに\R{単位}{たん|い}を\R{必}{かなら}ず\R{確}{たし}かめます。
\item \textbf{タイヤの\R{空気}{くう|き}とkPa}\par\noindent \R{車}{くるま}のタイヤの\R{空気}{くう|き}がおす\R{強}{つよ}さは、\qty{240}{kPa}（キロパスカル）くらいです。Pa（パスカル）も、Nと\R{同}{おな}じく\R{組}{く}み\R{立}{た}てた\R{単位}{たん|い}。kは1000\R{倍}{ばい}なので\mbox{\qty{240000}{Pa}}です。
\end{itemize}
\end{butsuri}

\begin{column}{1\,mはどうやって\R{決}{き}めた？}
1790\R{年}{ねん}ごろ、フランスで「\R{北極}{ほっ|きょく}から\R{赤道}{せき|どう}までの\R{長}{なが}さの1000\R{万分}{まん|ぶん}の1」を\qty{1}{m}と\R{決}{き}めました（2-8のコラムも\R{見}{み}てみよう）。\R{今}{いま}は、\R{光}{ひかり}が\R{真空}{しん|くう}の\R{中}{なか}を299792458\R{分}{ぶん}の1\R{秒}{びょう}に\R{進}{すす}む\R{長}{なが}さが\qty{1}{m}です。\R{光}{ひかり}の\R{速}{はや}さは、いつも\R{同}{おな}じだからです。\qty{1}{kg}も、2019\R{年}{ねん}から\R{自然}{し|ぜん}のきまりで\R{決}{き}めるようになりました。
\end{column}

\clearpage
%% ---------------- 章末問題 ----------------
\bkendproblems
\bkproblemgroup{\R{基本}{き|ほん}}{1}

\bkproblem \R{次}{つぎ}の\R{量}{りょう}を、（\quad）の\R{単位}{たん|い}で\R{表}{あらわ}しましょう。\par
（1）\nobreak \qty{4.2}{km}（m）\quad（2）\nobreak \qty{750}{mg}（g）\quad（3）\nobreak \qty{0.5}{L}（mL）\par
（4）\nobreak \qty{1.5}{m^2}（\unit{cm^2}）\quad（5）\nobreak \qty{0.2}{mm}（\unit{\micro m}）\quad（6）\nobreak \qty{40}{\micro m}（mm）
\bkpans{（1）\nobreak \qty{4200}{m}\quad（2）\nobreak \qty{0.75}{g}\quad（3）\nobreak \qty{500}{mL}\quad（4）\nobreak \qty{15000}{cm^2}\quad（5）\nobreak \qty{200}{\micro m}\quad（6）\nobreak \qty{0.04}{mm}}

\bkproblem 2\R{日}{にち}は\R{何}{なん}\R{分}{ぷん}ですか。\R{単位}{たん|い}つきの1をかける\R{式}{しき}で\R{書}{か}きましょう。
\bkpans{\qty{2}{d}×\qty{24}{h/d}×\qty{60}{min/h}＝2880\R{分}{ぷん}}

\bkproblem 1\R{年}{ねん}（365\R{日}{にち}）は\R{何}{なん}\R{秒}{びょう}ですか。365\R{日}{にち}を\qty{365}{d}として、\R{単位}{たん|い}つきの1を3つかける\R{式}{しき}を\R{書}{か}いて、\R{電卓}{でん|たく}で\R{計算}{けい|さん}しましょう。
\bkpans{\qty{365}{d}×\qty{24}{h/d}×\qty{60}{min/h}×\qty{60}{s/min}＝31536000\R{秒}{びょう}}

\bkproblem 3\R{分}{ぷん}20\R{秒}{びょう}で\qty{800}{m}\R{走}{はし}りました。\R{速}{はや}さは\R{何}{なん}m/sですか。
\bkpans{\qty{3}{min}×\qty{60}{s/min}＝\qty{180}{s}、\qty{180}{s}＋\qty{20}{s}＝\qty{200}{s}。$800\div200=4$ で\qty{4}{m/s}}

\bkproblem （1）\R{時速}{じ|そく}\qty{90}{km}は\R{秒速}{びょう|そく}\R{何}{なん}mですか。（2）\R{秒速}{びょう|そく}\qty{15}{m}は\R{時速}{じ|そく}\R{何}{なん}kmですか。
\bkpans{（1）\nobreak \qty{25}{m/s}\quad（2）\nobreak \qty{54}{km/h}}

\bkproblem \R{止}{と}まっていた\R{自転車}{じ|てん|しゃ}が、3\R{秒}{びょう}で\qty{6}{m/s}になりました。1\R{秒}{びょう}ごとに、\R{速}{はや}さは\R{何}{なん}m/sずつふえましたか（\R{同}{おな}じ\R{分}{ぶん}ずつふえたとする）。\R{単位}{たん|い}もつけて\R{答}{こた}えましょう。
\bkpans{$6\div3=2$ で\qty{2}{m/s^2}}

\bkproblem \R{次}{つぎ}の\R{式}{しき}のうち、\R{単位}{たん|い}が\R{合}{あ}わないものはどれですか。\par
\hspace*{1\zw}ア\enspace \R{速}{はや}さ＝\R{道}{みち}のり÷\R{時間}{じ|かん}\qquad イ\enspace \R{道}{みち}のり＝\R{速}{はや}さ÷\R{時間}{じ|かん}\par
\hspace*{1\zw}ウ\enspace \R{時間}{じ|かん}＝\R{道}{みち}のり÷\R{速}{はや}さ
\bkpans{イ（\unit{m/s}÷\unit{s}＝\unit{m/s^2}で、mにならない）}

\noindent\begin{minipage}[c]{0.58\linewidth}
\bkproblem \qty{2.4}{L}のジュースを、\qty{200}{mL}ずつコップに\R{分}{わ}けます。\R{単位}{たん|い}をそろえて、\R{何}{なん}ばいに\R{分}{わ}けられるか\R{求}{もと}めましょう。
\bkpans{\qty{2.4}{L}×\qty{1000}{mL/L}＝\qty{2400}{mL}。$2400\div200=12$ で12はい}
\end{minipage}\hfill
\begin{minipage}[c]{0.39\linewidth}\centering
\begin{tikzpicture}[x=1mm, y=1mm]
  % ジュースのびん（2.4 L）→ 200 mL ずつのコップ
  \bkbottle{(5,0)}{24}{0.85}
  \node[anchor=north, font=\sffamily\bfseries\small, inner sep=0pt] at (5,-1.8) {\qty{2.4}{L}};
  \draw[-{Stealth[length=2.2mm]}, bkInk!70, line width=0.9pt] (11.5,9) -- (18.5,9);
  \foreach \x in {24,33.5,43} {\bkcup{(\x,0)}{7.5}{10}{0.8}{0}}
  \node[font=\sffamily\bfseries\small, text=bkInk!70, inner sep=0pt] at (51.5,4.5) {…};
  \node[anchor=north, font=\sffamily\small, inner sep=0pt] at (33.5,-1.8) {\qty{200}{mL}ずつ};
\end{tikzpicture}
\end{minipage}

\clearpage
\bkproblemgroup{チャレンジ}{2}

\bkproblem \R{新幹線}{しん|かん|せん}の\R{速}{はや}さは\R{時速}{じ|そく}\qty{288}{km}、\R{陸上}{りく|じょう}100\,mの\R{選手}{せん|しゅ}の\R{速}{はや}さは\R{秒速}{びょう|そく}\qty{10}{m}です。\R{新幹線}{しん|かん|せん}は\R{選手}{せん|しゅ}の\R{何}{なん}\R{倍}{ばい}の\R{速}{はや}さですか。\R{単位}{たん|い}をそろえてくらべましょう。
\bkpans{$288\div3.6=80$ で\qty{80}{m/s}。$80\div10=8$ で8\R{倍}{ばい}}
\begin{bkcenter}
\begin{tikzpicture}[x=1mm, y=1mm]
  % 新幹線（ピクトグラム）
  \def\bkiikbody{(0,1.2) -- (0,10) -- (34,10) .. controls (45,10) and (52,6.4) .. (56,2.6) .. controls (56.9,1.7) and (56.4,1.2) .. (55.2,1.2) -- cycle}
  \draw[bkInk!60, line width=0.8pt] (-4,-1.2) -- (100,-1.2);
  \fill[white] \bkiikbody;
  \begin{scope}\clip \bkiikbody;
    \fill[bkBlue] (0,3.6) rectangle (60,5.2);
    \fill[bkBlue!75!black] (0,1.2) rectangle (60,2.4);
  \end{scope}
  \draw[bkInk!70, line width=0.7pt, line join=round] \bkiikbody;
  \foreach \x in {3,8,...,28} {\fill[bkInk!60, rounded corners=0.4pt] (\x,6.6) rectangle ++(3,2);}
  \fill[bkInk!75] (37,9.3) .. controls (41.5,9.3) and (44.5,8.2) .. (46.6,6.6) -- (38.4,6.6) -- cycle;
  \fill[bkInk!80] (6,0.2) circle (1.1) (10,0.2) circle (1.1) (42,0.2) circle (1.1) (46,0.2) circle (1.1);
  \node[anchor=south, font=\sffamily\bfseries\small, fill=bkBlue!12, draw=bkBlue!60, line width=0.5pt, rounded corners=1mm, inner sep=1.2mm] at (28,13) {\R{新幹線}{しん|かん|せん}　\R{時速}{じ|そく}\qty{288}{km}};
  % 陸上の選手
  \bkkid{(86,-1.2)}{17}{bkOrange}{bkHair}
  \draw[bkInk!50, line width=0.7pt, line cap=round] (77,5) -- (80.6,5) (76,8) -- (80.6,8) (77,11) -- (80.6,11);
  \node[anchor=south, font=\sffamily\bfseries\small, fill=bkOrange!12, draw=bkOrange!70, line width=0.5pt, rounded corners=1mm, inner sep=1.2mm] at (86,19) {\R{選手}{せん|しゅ}　\R{秒速}{びょう|そく}\qty{10}{m}};
\end{tikzpicture}
\end{bkcenter}

\bkproblem かみの\R{毛}{け}の\R{太}{ふと}さを\qty{80}{\micro m}とします。（1）これは\R{何}{なん}mmですか。（2）\qty{1}{mm}のはばに、かみの\R{毛}{け}はすき\R{間}{ま}なく\R{何}{なん}\R{本分}{ぼん|ぶん}\R{並}{なら}びますか。
\bkpans{（1）\nobreak \qty{0.08}{mm}\quad（2）\nobreak $1\div0.08=12.5$ で12.5\R{本分}{ほん|ぶん}}
\begin{bkcenter}
\begin{tikzpicture}[x=1mm, y=1mm]
  % 1 mm を 100 倍にした図（1 mm → 100 mm、かみの毛 80 μm → 8 mm）
  \fill[bkYellow!22] (0,0) rectangle (100,6);
  \draw[bkInk!70, line width=0.6pt] (0,0) rectangle (100,6);
  \foreach \k in {1,...,9} {\draw[bkInk!55, line width=0.4pt] (\k*10,6) -- (\k*10,3.8);}
  \draw[bkInk, line width=0.9pt] (0,6) -- (0,1.6) (100,6) -- (100,1.6);
  \node[anchor=north, font=\sffamily\footnotesize, inner sep=0pt] at (0,-1.6) {0};
  \node[anchor=north, font=\sffamily\footnotesize, inner sep=0pt] at (100,-1.6) {\qty{1}{mm}};
  \node[anchor=north, font=\sffamily\footnotesize, inner sep=0pt] at (50,-1.6) {（100\R{倍}{ばい}に\R{大}{おお}きくした\R{図}{ず}）};
  \foreach \k in {0,...,3} {\filldraw[fill=bkHair, draw=bkInk!35, line width=0.3pt] (\k*8,7.2) rectangle (\k*8+8,22);}
  \node[font=\sffamily\bfseries\large, text=bkInk!70, inner sep=0pt] at (42,14.6) {…};
  \node[font=\sffamily\bfseries\small, inner sep=0pt] at (68,14.6) {\R{何}{なん}\R{本分}{ぼん|ぶん}\R{並}{なら}ぶ？};
  \draw[{Stealth[length=1.6mm]}-{Stealth[length=1.6mm]}, bkRed, line width=0.6pt] (0,24.2) -- (8,24.2);
  \node[anchor=west, font=\sffamily\footnotesize, text=bkRed, inner sep=0pt] at (9.6,24.2) {かみの\R{毛}{け}1\R{本}{ぽん}\qty{80}{\micro m}};
\end{tikzpicture}
\end{bkcenter}


\begin{furikaeri}
\begin{checklist}
\item k・m・{\textmu} を\R{読}{よ}んで、\R{単位}{たん|い}を\R{書}{か}きかえられる。\hfill{\footnotesize\sffamily\color{bkTheme}→\pageref{bk:2-10:settou}ページ}
\item 「\R{単位}{たん|い}つきの1」で、\R{時間}{じ|かん}や\R{速}{はや}さを\R{換算}{かん|さん}できる。\hfill{\footnotesize\sffamily\color{bkTheme}→\pageref{bk:2-10:kansan}ページ}
\item \R{単位}{たん|い}を\R{計算}{けい|さん}して、\R{式}{しき}が\R{正}{ただ}しいかを\R{確}{たし}かめられる。\hfill{\footnotesize\sffamily\color{bkTheme}→\pageref{bk:2-10:keisan}ページ}
\end{checklist}
\end{furikaeri}

\bkfukidashi{sora}{\R{答}{こた}えに\R{単位}{たん|い}をつけると、まちがいに\R{自分}{じ|ぶん}で\R{気}{き}づけるんだね。\R{次}{つぎ}の\R{部}{ぶ}では、\R{文字}{も|じ}を\R{使}{つか}った\R{式}{しき}を\R{学}{まな}ぶよ。}
